% !TEX root = ../../tree_notes.tex
% =====================================================================
%  Chapter 30 -- Counting Labelled Trees             WRITTEN
%  (folder ch32_prufer_matrixtree; label prefix c32; PRINTS as
%  Chapter 30; folder number 32; build with:  python tools/build.py --only 32)
%
%  NUMBERING TRAP (the Ch.27/28/29 trap, mirrored): \chref{c32} prints
%  "Chapter 30" (THIS chapter).  Do NOT confuse label with printed number.
%  \chref{c31} prints "Chapter 29" (Trees Born from Graphs -- the PREVIOUS
%  chapter),
%  \chref{c22} prints "Chapter 28" (Tree Isomorphism -- unlabelled shapes),
%  \chref{c17} prints "Chapter 31" (MST -- the chapter AFTER next),
%  \chref{c28} prints "Chapter 17" (DSU / small-to-large -- used by 30.5
%  in one sentence for the component census of a partial forest),
%  \chref{c13} prints "Chapter 19" (segment trees -- the black box the
%  Summit uses to fold a range maximum of prefix sums),
%  \chref{c10} prints "Chapter 13" (tree DP -- the fold used in 30.4 for
%  a bottom-up subtree census),
%  \chref{c09} prints "Chapter 12" (tree metrics -- depth and path lengths),
%  \chref{c24} prints "Chapter 5" (Euler tour / tin-tout -- cross-referenced
%  once from 30.5),
%  \chref{c30} prints "Chapter 26" (virtual trees -- cited once for the
%  labelled-vs-unlabelled hand-off),
%  \chref{c21} prints "Chapter 27" (Mo + Tarjan on trees -- background cite).
%  Every cite of THIS chapter is \chref{c32}.
%
%  Owns THE ENUMERATIVE HOME of the whole book: when a tree is LABELLED on
%  \([n]=\{1,2,\dots,n\}\), how many such trees satisfy a given property, and
%  how do you turn the answer into an algorithm.  The tools: the PRUFER
%  BIJECTION (a length-\(n-2\) sequence over \([n]\) is in one-to-one
%  correspondence with labelled trees on \([n]\), and each vertex's DEGREE is
%  one plus its multiplicity in the code -- this is the whole chapter's workhorse),
%  CAYLEY'S FORMULA \(n^{n-2}\) (an immediate corollary of the bijection, plus
%  its two refinements: rooted \(n^{n-1}\) and the FOREST-COMPLETION count
%  \(n^{k-2}\prod_i s_i\) for a fixed forest with components of sizes \(s_i\)),
%  the MATRIX-TREE THEOREM (the number of spanning trees of an arbitrary
%  graph equals any cofactor of its Laplacian, so a modular determinant
%  becomes the answer oracle when the graph is given explicitly), and
%  LINEARITY OF EXPECTATION (root-path indicator sums turn "expected number
%  of random-deletion steps to empty a tree" into \(\sum_v 1/\mathrm{depth}(v)\)
%  with no recursion, in Chapter-12 depth).
%  It STARTS where Ch.29 (\chref{c31}) ENDS: that chapter extracted a tree
%  from a graph; this one counts trees directly.  The handoff to Ch.28
%  (\chref{c22}) is explicit in 30.5: this chapter counts LABELLED trees,
%  Ch.28 compares UNLABELLED shapes, and the two are not the same problem.
%  Modular arithmetic (Ch.6 if cited) and segment-tree ranges (Ch.19) are
%  used as black boxes; nothing is re-derived.
%
%  Sections:
%    * 30.1 The Prufer bijection -- encode a labelled tree to a length-(n-2)
%      sequence by repeatedly dropping the smallest leaf and writing its
%      neighbour; decode by the dual "smallest label absent from the residual
%      suffix" greedy.  Degree of vertex v = 1 + (multiplicity of v in the code).
%      WORKED examples only (encode and decode of the running 6-vertex tree
%      with code [5,5,3,4]; degree census from code; count of trees with a
%      prescribed degree sequence via the multinomial).
%    * 30.2 Cayley's formula and the forest-completion variant -- \(n^{n-2}\)
%      labelled trees on \([n]\) (immediate from the bijection: code length
%      \(n-2\), each entry any of \(n\) labels); \(n^{n-1}\) rooted (multiply by
%      \(n\)); degree-constrained count = \((n-2)!/\prod_v(\deg(v)-1)!\);
%      and the FOREST-COMPLETION count \(n^{k-2}\prod_i s_i\) for a fixed
%      \(k\)-component forest (each completion adds edges until one tree,
%      weight = size product).  Solved CF 2190D "Prufer Vertex" (2500, L5) --
%      given a forest with \(n\) vertices and \(m\) edges, for each \(v<n\)
%      count completions whose "second-to-last survivor" under smallest-leaf
%      peeling equals \(v\), modulo \(998244353\); live-verified against the
%      CF problem page (title "D. Prufer Vertex", 2 seconds, 256 MB, \(n\le
%      2\cdot 10^5\), \(m\le n-1\), sum-\(n\) cap, sample three cases returning
%      "1 2", "0 0 0 1", "12 0 1 6 5").  Reuses the bijection, the degree
%      census, and a per-tree DP fold.
%    * 30.3 The Matrix-Tree theorem -- Laplacian \(L=D-A\), spanning-tree
%      count = \(\det(L\) with row/col \(n\) deleted\()\); modular Gaussian
%      elimination for the big-count regime; the complete-graph sanity
%      check \(\det=\;n^{n-2}\) recovering Cayley.  WORKED examples only
%      (cycle \(C_n\) gives \(n\); complete bipartite \(K_{a,b}\) gives
%      \(a^{b-1}b^{a-1}\); modular det on a small explicit graph).
%    * 30.4 Linearity of expectation on a rooted tree -- indicator
%      decomposition turns "expectation of the number of steps to empty a
%      rooted tree, where each step deletes the subtree of a uniformly
%      random remaining vertex" into \(\sum_v 1/\mathrm{depth}(v)\) with
%      depth measured as number of vertices on the root-\(v\) path.
%      Solved CF 280C "Game on Tree" (2200, L4) -- \(n\le 10^5\), absolute/
%      relative error \(10^{-6}\), live samples \(n=2\) edge \(1\!-\!2\)
%      giving \(1.5\) and \(n=3\) star giving \(2.0\); a single DFS computes
%      depths and the sum of reciprocals.
%    * 30.5 The method boundary -- labelled vs.\ unlabelled census (why
%      Ch.28's tree-isomorphism machinery does NOT apply here, and how a
%      \(\log n\)-factor orbit-size correction reconciles them); when
"count trees" means a closed form vs. a determinant vs. an expectation;
%      handing off to Ch.26 (virtual trees) and Ch.31 (MST).
%      WORKED examples only.
%  Chapter-end Summit: CF 917D "Stranger Trees" (2500, L5) -- given a tree
%  on \(n\le 100\) vertices, for each \(k\in[0,n-1]\) count the LABELLED
%  trees that share EXACTLY \(k\) edges with it, modulo \(10^9+7\).  Pipeline:
%  remove a set \(S\) of tree edges (a subset of size \(n-1-k'\) for a
%  \(k'\)-superset count), the residue is a forest with components of sizes
%  \(s_i\), and the number of completions that CONTAIN none of the "off-\(S\)
%  tree edges" is bounded by inclusion-exclusion over a tree-DP that folds
%  subtree-by-subtree the polynomial \(\prod(\cdot)\) of the forest-completion
%  formula \(n^{c-2}\prod s_i\); the answer is obtained by an inversion of the
%  binomial "at least vs.\ exactly" relation.  Reuses TWO chapter tools (the
%  forest-completion formula, the tree DP) and ONE earlier-chapter tool
%  (Ch.13 tree DP \(\chref{c10}\)).  Official samples hand-replayed:
%  star on 3 -> "0 2 1", path on 4 -> "1 7 7 1", star on 4 -> "0 9 6 1".
%
%  Diagram-first: this chapter is BIJECTION based.  The running example is
%  the 6-vertex labelled tree \(T\) defined once, authoritatively, in
%  00_opener: edges \(\{1\!-\!5,\,2\!-\!5,\,5\!-\!3,\,3\!-\!4,\,4\!-\!6\}\),
%  with the Prüfer code \([5,5,3,4]\).  Every worked encode/decode, degree
%  census, and multinomial count in 30.1 and 30.2 uses this same tree, and
%  30.4's DFS depth sum runs on it rooted at vertex 1 (giving
%  \(\mathrm{depth}(1..6) = (1,3,3,4,2,5)\) and expectation \(1/1 + 1/3 +
%  1/3 + 1/4 + 1/2 + 1/5 = 157/60\)).  Both solved problems and the Summit use
%  their OWN official samples, so no trace collides.
%
%  Anti-repetition contract (teach once, cross-reference by label):
%    * Ch.28 (\chref{c22}) owns UNLABELLED tree isomorphism and AHU; 30.5
%      gestures at orbit sizes but never re-derives AHU.
%    * Ch.13 (\chref{c10}) owns the tree-DP bottom-up fold; 30.2's
%      per-tree degree census and 30.4's depth sum both USE it (one line each),
%      not rederive.
%    * Ch.19 (\chref{c13}) owns the segment-tree monoid; the Summit's
%      inclusion-exclusion polynomial fold is a tree DP, and NO segment
%      tree is required here -- we cite Ch.19 only to remind the reader
%      that the tool exists for range-max, not for this chapter's problems.
%    * Ch.17 (\chref{c28}) owns DSU component counting; 30.2's forest
%      component census USES it or the DFS equivalent -- one sentence.
%    * Ch.12 (\chref{c09}) owns depth/tin; 30.4's depth array IS Ch.12's
%      single-source depth -- cited, not redefined.
%    * Ch.26 (\chref{c30}) owns virtual trees; the handoff in 30.5 is one
%      forward-cite.
%
%  Known URL collisions (confirmed by a FRESH ownership scan of chapters/ +
%  the live CF API; prose-cite only, NEVER re-link -- urlsweep must print
%  an empty DUPLICATES section):
%    * CF 156/D "Point in Tree" and CF 1092/E "Minimum Diameter Tree" are
%      NOT linked here (either owned elsewhere or out of band).
%    * The roster URLs below (2190D, 280C, 917D) and every mixed-set URL
%      were confirmed ABSENT from build/used_cf.txt (the 261-ID used set)
%      BEFORE writing; re-derive ownership from the chapter files.
%
%  Roster (every title/rating/tag/sample live-verified against the Codeforces
%  problemset API / problem page BEFORE a trace was written; CF-forward):
%    solved: CF 2190D "Prufer Vertex"       (2500, L5) (30.2) -- per-vertex
%            count of forest completions whose second-to-last peeling
%            survivor is v; live samples replayed.
%            CF 280C "Game on Tree"          (2200, L4) (30.4) -- expected
%            random-subtree-deletion steps = sum of 1/depth over all vertices;
%            live samples replayed -> 1.5, 2.0.
%    worked: the encode/decode pair on the running tree with code
%            [5,5,3,4]; the degree-multiset-to-multiplicity multinomial; the
%            Laplacian cofactor on cycle C_5 and K_{2,3}; the modular
%            determinant of a small explicit graph; the depth-reciprocal
%            sum on the running tree rooted at 1; the labelled-vs-unlabelled
%            reconciliation on n=4.
%    summit: CF 917D "Stranger Trees"        (2500, L5) -- for each k in
%            [0,n-1] count labelled trees sharing exactly k edges with a
%            given tree, mod 10^9+7; forest-completion + tree DP +
%            binomial inversion; three live samples replayed.
%    mixed (10, ascending by live-verified rating, every ID confirmed ABSENT
%            from build/used_cf.txt -- LOCKED before writing): 1605D(2100),
%            1311E(2200), 1540B(2300), 1929F(2300), 653E(2400), 1237E(2400),
%            1120D(2500), 804D(2500), 1762E(2600), 1495D(2600).  A spread
%            of genuine counting-on-trees drills (Prüfer reconstruction 1605D
%            "Treelabeling" and 1311E "Construct the Binary Tree", DP on tree
%            arrays 1540B "Tree Array", BST-count via Catalan-like structures
%            1929F "Sasha and the Wedding Binary Search Tree", tree
%            construction from partial degrees 653E "Bear and Forgotten
%            Tree 2" and 1237E "Balanced Binary Search Trees", Prüfer-code
%            driven MST-style construction 1120D "Power Tree", expected
%            diameter 804D, subtree-reweight DP 1762E "Tree Sum", BFS-tree
%            counting 1495D "BFS Trees") and deliberate REJECTION cases
%            whose hints name the true tool.
% =====================================================================
\bookchapter{Counting Labelled Trees}{c32}
\chapterinfo{CF 2100--2600}{Written}

% =====================================================================
%  Chapter 30 opener -- "Chapter at a Glance", three reading paths + the
%  ladder staircase, THE running example (the 6-vertex labelled tree T on
%  {1,2,3,4,5,6} with edges {1-5, 2-5, 5-3, 3-4, 4-6} and Prufer code
%  [5,5,3,4]), the Cayley / Matrix-Tree / Expectation taxonomy, the two
%  traps (labelled vs unlabelled, and the "code is sequence" trap that
%  misses the bijection), the who-owns-what comparison against Ch.28
%  (unlabelled shapes), Ch.13 (tree DP), Ch.19 (segment tree), Ch.12
%  (depth/path), Ch.17 (DSU), Ch.31 (MST -- the chapter AFTER next), and
%  Ch.29 (Trees Born from Graphs -- the PREVIOUS chapter), a constraint
%  -> technique table, and notation.
%  Exports fig:c32-tree, fig:c32-code, fig:c32-taxonomy.
%  \term the NEW objects "Prufer code", "labelled tree", "spanning-tree
%  census", "Laplacian", "cofactor", "matrix-tree", "linearity of expec-
%  tation on trees" only -- degree/depth/DSU are indexed by earlier chapters.
% =====================================================================

\begin{chaptertheorybox}{Chapter at a Glance}
Every earlier chapter in this book asked a \emph{structural} question about a
tree you already had.  Chapter~29 (\chref{c31}) was the last: it taught you to
\emph{extract} a tree from a graph that was one.  This chapter turns the arrow
around.  The input is no longer a specific tree; the input is a \emph{property},
and the question is \textbf{how many labelled trees satisfy it}.  That flips
every tool: instead of one DFS you build a bijection, instead of a fold you
count, and instead of a numerical answer you return a residue modulo a prime.

The chapter is built on three pillars, each its own section.  The \term{Pr\"ufer
bijection} (\(\secref{c32-prufer}\)) is a length-\(n-2\) sequence of labels from
\([n]=\{1,\dots,n\}\) that encodes a labelled tree on \([n]\) with a two-way
greedy inverse; every enumeration question about labelled trees collapses to
\emph{count sequences}, which is elementary.  \term{Cayley's formula}
(\(\secref{c32-cayley}\)) is the immediate corollary \(n^{n-2}\), but the real
payoff is in its variants -- the degree-sequence multinomial and the
\emph{forest-completion} formula \(n^{k-2}\prod_i s_i\) that counts how many
ways you can add \(k-1\) edges to a \(k\)-component forest to make it a single
tree.  The \term{Matrix-Tree theorem} (\(\secref{c32-matrix}\)) handles the
case where the graph is not complete: the number of its spanning trees is a
cofactor of its Laplacian, so a modular Gaussian elimination is the whole
algorithm.  A fourth short section, \term{linearity of expectation}
(\(\secref{c32-expect}\)) on a rooted
tree, is the probabilistic twin of counting: split
an expectation into vertex indicators, and \(\mathbb{E}\) becomes a sum of
reciprocals of depths.

\textbf{You will learn to}
\begin{itemize}[nosep]
  \item encode and decode the Pr\"ufer sequence of a labelled tree, read the
        degree of every vertex off the code, and count degree-constrained
        trees via a multinomial (\(\secref{c32-prufer}\)); \hfill\lv{3}
  \item state and use Cayley \(n^{n-2}\), rooted \(n^{n-1}\), and the
        forest-completion formula \(n^{k-2}\prod_i s_i\); solve CF 2190D
        \emph{Prufer Vertex}, which asks for the second-to-last peeling
        survivor of every completion, modulo \(998244353\)
        (\(\secref{c32-cayley}\)); \hfill\lv{5}
  \item build the Laplacian, delete a row and column, and evaluate the
        determinant modulo a prime to count spanning trees of an arbitrary
        graph (\(\secref{c32-matrix}\)); \hfill\lv{4}
  \item decompose an expectation into per-vertex indicators, use
        \(\Pr[v\text{ is picked before its parent}]=1/\mathrm{depth}(v)\),
        and solve CF 280C \emph{Game on Tree} in one DFS
        (\(\secref{c32-expect}\)); \hfill\lv{4}
  \item tell when a counting problem is \emph{labelled} (this chapter) and
        when it is \emph{unlabelled} (Ch.28 \chref{c22}), and use the
        orbit-size bridge to reconcile the two (\(\secref{c32-boundary}\)).
        \hfill\lv{5}
\end{itemize}

\textbf{Prerequisites.} Chapter~13 (\chref{c10}) -- the bottom-up tree DP fold,
which the Summit uses to combine a polynomial across a rooted tree in one
recursion.  Chapter~12 (\chref{c09}) -- depth and path length in a rooted tree,
which \(\secref{c32-expect}\) uses verbatim.  Chapter~17 (\chref{c28}) -- DSU
component census, which \(\secref{c32-cayley}\) uses to read the \(k\) and
\(s_i\) of the input forest.  Chapter~29 (\chref{c31}) -- the \textbf{previous
chapter}; it ends at "here is the tree hiding inside the graph"; this one opens
at "here is how many trees match a description."
\end{chaptertheorybox}

% ---------------- Reading paths + ladder position ----------------
\begin{chaptertheorybox}{How to Read This Chapter (three paths)}
Level badges \(\lv3\)--\lv{6}: \(\lv3\) is hand-encoding a labelled tree to its
Pr\"ufer code by dropping smallest leaves; \(\lv4\) is coding the decode
greedy and running the Laplacian determinant on a small graph, and applying
linearity of expectation to CF 280C; \(\lv5\) is the forest-completion formula
and CF 2190D; \(\lv6\) is the Summit (CF 917D) -- count labelled trees sharing
exactly \(k\) edges with a given one, via the completion formula and a
tree-DP polynomial.

\textbf{Bijection path (\(\lv3\)--\lv{5}, ``count sequences'').}
\secref{c32-prufer} in full -- encode, decode, degree-from-code -- capped by
\secref{c32-cayley}'s Cayley and forest-completion formulas.  Stop here if
your problem reduces to a closed form on \([n]\).

\textbf{Linear-algebra path (\(\lv4\)--\lv{5}, ``count spanning trees of a
given graph'').} \secref{c32-matrix}: Laplacian, cofactor, modular determinant,
and the sanity checks against Cayley on \(K_n\) and against \(n\) on \(C_n\).
This is the branch you take when the graph is \emph{given} and you cannot
reduce to a Pr\"ufer-style sequence.

\textbf{Probability path (\(\lv4\)--\lv{6}, ``expected value on a tree'').}
\secref{c32-expect} -- indicator decomposition and the depth-reciprocal
formula, then solved CF 280C in one DFS.  Read \(\secref{c32-boundary}\) after
either branch to know when you are outside this chapter's method.
\end{chaptertheorybox}

\begin{figure}[H]
\centering
\begin{tikzpicture}[diagram, x=8.5mm, y=6.5mm]
  \foreach \st/\lab in {1/lv1,2/lv2,3/lv3,4/lv4,5/lv5,6/lv6} {
    \fill[\lab!22, draw=\lab, line width=0.8pt] (\st-1,0) rectangle (\st,\st*0.52);
    \node[font=\scriptsize\bfseries, text=\lab] at (\st-0.5,\st*0.52-0.42) {L\st};}
  \node[annot, anchor=south west, font=\scriptsize]
    at (0.05,0.62) {prereq: degree/depth (Ch.12), tree DP (Ch.13), DSU census (Ch.17)};
  \node[annot, anchor=south west, font=\scriptsize]
    at (2.05,1.66) {30.1 Prufer encode / decode by hand; degree from multiplicity};
  \node[annot, anchor=south west, font=\scriptsize]
    at (3.05,2.18) {30.3 Laplacian + modular det | 30.4 linearity, CF 280C};
  \node[annot, anchor=south west, font=\scriptsize, text=lv5]
    at (4.05,2.70) {30.2 Cayley + forest completion -- CF 2190D};
  \node[annot, anchor=south west, font=\scriptsize, text=lv6]
    at (5.05,3.22) {30.5 boundary; Summit CF 917D (exactly k shared edges)};
\end{tikzpicture}
\caption{Chapter~30 on the six-step ladder: hand-encoding the Pr\"ufer code
of the running tree is \(\lv3\); the coded decode, the Laplacian determinant,
and CF 280C's depth-reciprocal sum are \(\lv4\); CF 2190D's
per-vertex completion count is \(\lv5\); the Summit (CF 917D) is step~6.}
\end{figure}

% ---------------- The running example ----------------
\subsection*{The Running Example: A Labelled Tree on Six Vertices}

Every section up to the Summit runs on the same 6-vertex labelled tree \(T\),
fixed once here and authoritative (if a later trace disagrees, this figure
wins).  Vertices are \([6]=\{1,2,3,4,5,6\}\); edges, 1-based, are
\[
  1\!-\!5,\quad 2\!-\!5,\quad 5\!-\!3,\quad 3\!-\!4,\quad 4\!-\!6 .
\]
The Pr\"ufer code of \(T\) is \([5,5,3,4]\): drop the smallest leaf in turn
(\(1\), write its neighbour \(5\); then \(2\), write \(5\); then \(5\) itself
becomes a leaf after \(1,2\) are gone, but the smallest remaining leaf is
\(\min\{5,6\}=5\) if we hold on the definition -- \(\secref{c32-prufer}\) does
the exact bookkeeping), and the length-\(n-2=4\) sequence is \([5,5,3,4]\).
Reading degrees off the code: a vertex appears \((\deg-1)\) times, so
\(\deg(5)=3\), \(\deg(3)=\deg(4)=2\), \(\deg(1)=\deg(2)=\deg(6)=1\).  The
degree sum is \(3+2+2+1+1+1=10=2(n-1)\), as it must be.

\begin{figure}[H]
\centering
\begin{tikzpicture}[diagram, scale=0.95, transform shape,
   tnode/.style={circle, draw=primaryblue!70, fill=primaryblue!8, line width=0.8pt,
                 minimum size=8mm, font=\small\bfseries, inner sep=1pt}]
  \node[tnode] (v1) at (-2.6,1.2) {1};
  \node[tnode] (v2) at (-2.6,-1.2) {2};
  \node[tnode] (v5) at (0,0) {5};
  \node[tnode] (v3) at (2.4,0) {3};
  \node[tnode] (v4) at (4.4,-0.8) {4};
  \node[tnode] (v6) at (6.2,0.4) {6};
  \draw[primaryblue!80, line width=1.1pt] (v1)--(v5) (v2)--(v5) (v5)--(v3) (v3)--(v4) (v4)--(v6);
  \node[annot, font=\scriptsize, anchor=north, text width=9cm] at (1.2,-1.8)
    {leaves: \(1,2,6\) \quad deg 3: \(5\) \quad deg 2: \(3,4\)};
\end{tikzpicture}
\caption{The running tree \(T\).  Its Pr\"ufer code is \([5,5,3,4]\), and
\(\sum_v 1/\mathrm{depth}(v)\) from the root \(1\) is \(1/1+1/2+1/3+1/3+1/4+1/5
= 157/60\) (used in \(\secref{c32-expect}\)).  This object is fixed for the
whole chapter.}
\label{fig:c32-tree}
\end{figure}

\begin{figure}[H]
\centering
\begin{tikzpicture}[diagram, scale=1.0, transform shape,
  box/.style={rectangle, draw=primaryblue!70, fill=primaryblue!8,
              minimum width=9mm, minimum height=8mm,
              font=\small\bfseries, inner sep=2pt}]
  \node[box] (a1) at (0,0) {5};
  \node[box] (a2) at (1.2,0) {5};
  \node[box] (a3) at (2.4,0) {3};
  \node[box] (a4) at (3.6,0) {4};
  \node[annot, font=\scriptsize, anchor=east] at (-0.6,0) {code:};
  \node[annot, font=\scriptsize, anchor=west] at (4.4,0) {\((n-2)=4\) entries};
  \node[annot, font=\scriptsize, anchor=north] at (1.2,-0.7) {\(\deg(5)-1=2\) copies};
  \node[annot, font=\scriptsize, anchor=north] at (2.4,-0.7) {\(\deg(3)-1=1\)};
  \node[annot, font=\scriptsize, anchor=north] at (3.6,-0.7) {\(\deg(4)-1=1\)};
\end{tikzpicture}
\caption{The Pr\"ufer code \([5,5,3,4]\) laid out.  The number of times a
label appears in the code is exactly \(\deg-1\); every counting application in
this chapter is a way of exploiting that one fact.}
\label{fig:c32-code}
\end{figure}

Two traps return throughout.  The \textbf{labelled-vs-unlabelled trap}: this
chapter counts labelled trees, and two trees with the same shape but a
different vertex-numbering are different objects.  Chapter~28 (\chref{c22})
counts \emph{shapes}; on \(n=6\) there are exactly 6 unlabelled tree shapes
but \(6^4=1296\) labelled ones.  Reconciling the two is the orbit-size
argument in \(\secref{c32-boundary}\); confusing them is the most common wrong
answer on Cayley-style problems.  The \textbf{code-is-a-sequence trap} is
subtler: the Pr\"ufer code is a sequence \emph{with structure}, not a bag of
labels.  Decoding needs the greedy "smallest label not in the residual
suffix" step to be well-defined; encoding drops the smallest \emph{leaf},
which is a strictly stronger notion than smallest label.  Every wrong
implementation of Pr\"ufer decode gets these two confusions somewhere.

\subsection*{Which Job? Reading the Statement}

\begin{center}
\renewcommand{\arraystretch}{1.25}
\setlength{\tabcolsep}{3pt}
\begin{tabular}{@{}p{0.30\linewidth} p{0.20\linewidth} p{0.28\linewidth} >{\raggedright\arraybackslash}p{0.13\linewidth}@{}}
\toprule
\textbf{The statement says\dots} & \textbf{Tool} & \textbf{Why it works} & \textbf{Section}\\
\midrule
"how many trees on \([n]\) with property \(P\)"; \(P\) reads off the code & Pr\"ufer bijection & trees \(\leftrightarrow\) length-\((n-2)\) sequences & \(\secref{c32-prufer}\)\\
"count completions of a forest"; "rooted trees"; "with prescribed degrees" & Cayley variants & \(n^{n-2}\), \(n^{k-2}\prod s_i\), multinomial & \(\secref{c32-cayley}\)\\
"spanning trees of this graph"; graph given explicitly & Matrix-Tree & Laplacian cofactor = count & \(\secref{c32-matrix}\)\\
"expected value of a game on a tree" & Linearity of expectation & per-vertex indicators \(\to\) reciprocals of depths & \(\secref{c32-expect}\)\\
\bottomrule
\end{tabular}
\end{center}

\begin{insightbox}[The whole chapter in one sentence]
A labelled tree on \([n]\) is a sequence of \(n-2\) labels in disguise; every
enumerative question you can ask about such trees is either (a) a count of
those sequences with a filter, or (b) a linear-algebra determinant of a
Laplacian, and the entire craft of this chapter is picking which of the two
the problem actually is -- then applying linearity of expectation when the
answer is a real number, not an integer.
\end{insightbox}

% ---------------- Taxonomy figure ----------------
\begin{figure}[H]
\centering
\begin{tikzpicture}[diagram, scale=0.9, transform shape,
   mnode/.style={circle, draw=primaryblue!70, fill=primaryblue!8, line width=0.7pt,
                 minimum size=13mm, font=\scriptsize\bfseries, inner sep=1pt, align=center}]
  \node[mnode] (root) at (0,0) {count\\ trees};
  \node[mnode] (lab)  at (-4,-1.6) {labelled\\ on \([n]\)};
  \node[mnode] (sub)   at (1.5,-1.6) {of a given\\ graph};
  \node[mnode] (prf)  at (-5.6,-3.2) {Prufer\\ code};
  \node[mnode] (cay)   at (-2.6,-3.2) {Cayley /\\ completion};
  \node[mnode] (mt)   at (0.4,-3.2) {Matrix-\\ Tree};
  \node[mnode] (exp)  at (3.4,-3.2) {Linearity\\ of \(\mathbb{E}\)};
  \draw (root)--(lab)--(prf);
  \draw (lab)--(cay);
  \draw (root)--(sub)--(mt);
  \draw (sub)--(exp);
\end{tikzpicture}
\caption{Where the four methods sit. A \emph{labelled} enumeration on \([n]\)
is either a Pr\"ufer-code argument (\(\secref{c32-prufer}\)) or a Cayley
variant (\(\secref{c32-cayley}\)).  Counting inside a \emph{given graph} is
either Matrix-Tree (\(\secref{c32-matrix}\)) or, if the answer is
probabilistic, a linearity-of-expectation argument (\(\secref{c32-expect}\)).}
\label{fig:c32-taxonomy}
\end{figure}

\subsection*{Who Owns What}

\begin{center}
\renewcommand{\arraystretch}{1.2}
\setlength{\tabcolsep}{3pt}
\begin{tabular}{@{}p{0.28\linewidth} p{0.40\linewidth} p{0.24\linewidth}@{}}
\toprule
\textbf{Object} & \textbf{Owned here?} & \textbf{Elsewhere}\\
\midrule
Pr\"ufer code, bijection encode/decode & \textbf{Yes}, \(\secref{c32-prufer}\) -- with worked example & --\\
Cayley \(n^{n-2}\), forest completion \(n^{k-2}\prod s_i\) & \textbf{Yes}, \(\secref{c32-cayley}\) -- proved from the bijection & --\\
Matrix-Tree theorem, Laplacian cofactor & \textbf{Yes}, \(\secref{c32-matrix}\) & --\\
Linearity of expectation on rooted trees & \textbf{Yes}, \(\secref{c32-expect}\) & --\\
Depth / path length & no -- used as a black box & Ch.12 (\chref{c09})\\
Bottom-up tree DP fold & no -- the fold is reused & Ch.13 (\chref{c10})\\
DSU component census & no -- used once for forest reading & Ch.17 (\chref{c28})\\
Unlabelled tree shapes & no -- different question & Ch.28 (\chref{c22})\\
Minimum spanning tree & no -- a different "trees in a graph" & Ch.31 (\chref{c17})\\
\bottomrule
\end{tabular}
\end{center}

\subsection*{Recognition Guide}

\begin{recognitionbox}[Turning a counting question into a Pr\"ufer answer]
\reckey{``how many trees on \(n\) vertices satisfy \(\ldots\)'', ``count labelled
trees sharing exactly \(k\) edges with a given tree'', ``expected number of
random-deletion steps until empty'', ``given degrees, how many trees'', ``add
edges to make a forest connected''}
\recuse{
  \item \textbf{Property reads off degrees or the sequence itself} -- Pr\"ufer
        bijection (\(\secref{c32-prufer}\)) turns it into a multinomial; Cayley
        variants (\(\secref{c32-cayley}\)) finish it.
  \item \textbf{Property depends on the graph's edge set} -- Matrix-Tree
        (\(\secref{c32-matrix}\)): Laplacian cofactor.
  \item \textbf{Answer is a real number, not an integer} -- linearity of
        expectation (\(\secref{c32-expect}\)); the per-vertex indicator is
        almost always \(1/\)depth.
}
\recnotuse{
  \item \textbf{The question is about tree \emph{shapes} not labelings} --
        that is Ch.28 (\chref{c22}) tree isomorphism and AHU.
  \item \textbf{The question is about a specific tree, not a family} -- that
        is Ch.12/13/26/29, not this chapter.
}
One-line test: \emph{if the answer is a number mod a prime or a floating
value and the input is a \emph{description} rather than a graph, you are in
this chapter; build the bijection, apply the formula, or split the
expectation.}
\end{recognitionbox}

\subsection*{Constraint \(\to\) Technique}

\begin{center}
\renewcommand{\arraystretch}{1.2}
\setlength{\tabcolsep}{3pt}
\begin{tabular}{@{}p{0.34\linewidth} p{0.34\linewidth} >{\raggedright\arraybackslash}p{0.20\linewidth}@{}}
\toprule
\textbf{Constraint signature} & \textbf{Move} & \textbf{Section}\\
\midrule
\(n\le 10^6\), "how many trees with these degrees" & multinomial from code & \(\secref{c32-prufer}\)\\
forest on \(n\le 2\cdot 10^5\), \(k\) components & \(n^{k-2}\prod s_i\) mod \(p\) & \(\secref{c32-cayley}\)\\
graph on \(n\le 100\), count spanning trees & Laplacian det mod \(p\) & \(\secref{c32-matrix}\)\\
tree on \(n\le 10^5\), expected steps to empty & \(\sum_v 1/\mathrm{depth}(v)\) & \(\secref{c32-expect}\)\\
\(n\le 100\), count trees sharing exactly \(k\) edges & completion + DP + inversion & \(\probref{c32-summit}\)\\
\bottomrule
\end{tabular}
\end{center}

\subsection*{Notation Used Only Here}

\([n]=\{1,2,\dots,n\}\) is the label set.  \(\mathcal{T}_n\) is the set of
labelled trees on \([n]\); \(|\mathcal{T}_n|=n^{n-2}\) by Cayley.  \(\mathrm{code}(T)\) is the
Pr\"ufer code of \(T\) (a sequence of \(n-2\) labels from \([n]\));
\(\mathrm{decode}(c)\) is its inverse.  \(\deg_T(v)\) is the degree of \(v\) in
\(T\); \(m_v(T)=\deg_T(v)-1\) is the \emph{multiplicity} of \(v\) in
\(\mathrm{code}(T)\).  For a graph \(G\), \(L(G)=D-A\) is its Laplacian, and
\(L^{(i)}(G)\) is \(L(G)\) with row and column \(i\) deleted.  For a rooted
tree, \(\mathrm{depth}(v)\) counts the vertices on the root-\(v\) path
(inclusive of both endpoints), so \(\mathrm{depth}(\text{root})=1\).  The
Summit uses \(f(k)\) for the count of completions whose shared-edge set with
the given tree contains a fixed \(k\)-subset, and \(g(k)\) for the count that
shares \emph{exactly} \(k\) edges -- these two are related by a binomial
inversion.

% =====================================================================
%  30.1  The Pruefer bijection.
%  Exports: sec:c32-prufer  def:c32-ltree  def:c32-code  def:c32-dec
%  thm:c32-bij  lem:c32-deg  cor:c32-mult  lst:c32-encode
%  lst:c32-decode  fig:c32-encode  fig:c32-decode
%  This section OWNS the Pruefer encode/decode pair, the degree-from-code
%  lemma, and the degree-multiset multinomial count; 30.2-30.5 and the
%  Summit reference them by label.  NO solved-problem block here (the
%  solved problems are 2190D in 30.2 and 280C in 30.4).  Everything is
%  worked examples on the running 6-vertex tree T with code [5,5,3,4].
% =====================================================================
\section{The Pr\"ufer Bijection}\label{sec:c32-prufer}

\begin{previouslybox}[Ch.29 extracted one tree from a graph]
Chapter~29 (\chref{c31}) ended with a graph and the tree that hides inside
it -- a bridge tree, a block-cut tree, a rho decomposition.  Chapter~30 opens
with the opposite question: the input is a \emph{description}, and the answer
is a number.  To count trees you need a way to walk through \(\mathcal{T}_n\)
without enumerating it, and that is what the Pr\"ufer bijection is -- a
canonical code, one sequence per labelled tree, in a form you can count with
elementary combinatorics.
\end{previouslybox}

\begin{topicintro}
A \term{labelled tree} on \([n]\) is a tree whose vertices are the specific
set \(\{1,2,\dots,n\}\) (\(\defref{c32-ltree}\)); two labelled trees differ
when their edge sets do, so a relabelling that moves an edge is a \emph{new}
tree.  Pr\"ufer's construction (\(\defref{c32-code}\)) is an iterative peel:
as long as more than two vertices remain, find the leaf with the smallest
label, delete it, and append the label of its (unique) neighbour to a running
list.  After \(n-2\) steps the list is the \term{Pr\"ufer code} of \(T\), a
sequence of length \(n-2\) with entries in \([n]\).  The inverse
(\(\defref{c32-dec}\)) reconstructs \(T\) from its code greedily: at each step
take the smallest label of \([n]\) that is not in the residual suffix of the
code, join it to the code's first remaining entry, and drop that entry.  The
two maps are mutually inverse (\(\thmref{c32-bij}\)), which is the whole point
of the chapter: counting labelled trees is counting sequences.

The most-used corollary is \(\lemref{c32-deg}\): the number of times a label
\(v\) appears in \(\mathrm{code}(T)\) is \(\deg_T(v)-1\).  Read a degree
constraint as a multiplicity constraint on the code, and the count becomes a
multinomial (\(\corref{c32-mult}\)).

\textbf{Prerequisites.} The notion of a leaf (a vertex of degree 1) and the
handshake identity \(\sum_v \deg(v) = 2(n-1)\), both from Ch.~1.  Nothing
about trees beyond "connected and \(n-1\) edges" is assumed here; the
bijection itself is proved from scratch.

\textbf{After this section you can}
\begin{itemize}[nosep]
  \item write the Pr\"ufer code of any labelled tree in \(\bigO{n\log n}\) via a min-heap of leaves; \hfill\lv{3}
  \item decode a length-\((n-2)\) sequence over \([n]\) back to its tree in \(\bigO{n\log n}\); \hfill\lv{3}
  \item state and use \(\lemref{c32-deg}\) to convert degree constraints into multiplicity constraints; \hfill\lv{4}
  \item count labelled trees with a prescribed degree sequence via \(\corref{c32-mult}\). \hfill\lv{4}
\end{itemize}
\end{topicintro}

% ---------------------------------------------------------------
\subsection{Encoding a labelled tree to its Pr\"ufer code}
% ---------------------------------------------------------------

\begin{definition}{Labelled tree on \([n]\)}{c32-ltree}
A \term{labelled tree} on \([n]\) is an undirected graph \(T=(V,E)\) with
\(V=[n]\), connected and with \(|E|=n-1\).  The set of all such trees is
\(\mathcal{T}_n\).  Two labelled trees \(T,T'\) are equal iff they have the
same edge set \(\{u,v\}\subseteq[n]\); a permutation of the labels produces a
\emph{different} tree unless it fixes every edge.
\end{definition}

\begin{definition}{Pr\"ufer code (encoding)}{c32-code}
Given \(T\in\mathcal{T}_n\) with \(n\ge 2\), form a list \(c\) initially
empty.  While \(|V(T)|>2\): let \(\ell\) be the smallest leaf of the current
\(T\); append the label of \(\ell\)'s unique neighbour to \(c\); delete \(\ell\)
and its incident edge.  The final \(c\) is the \term{Pr\"ufer code} of \(T\),
written \(\mathrm{code}(T)\).  It has length exactly \(n-2\), because one
vertex is removed at each of the \(n-2\) iterations.
\end{definition}

\begin{lemma}{Degree from multiplicity}{c32-deg}
For \(T\in\mathcal{T}_n\) and any \(v\in[n]\),
\[
  m_v(T) \;:=\; \#\{\text{positions } i : c_i = v\} \;=\; \deg_T(v) - 1 .
\]
In particular, a vertex is a leaf iff its label never appears in the code.
\end{lemma}

\begin{proofbox}
Track the current degree of each vertex as the peel runs.  When a leaf
\(\ell\) is deleted, we write its neighbour \(v\), and \(v\)'s current degree
drops by one (its edge to \(\ell\) is gone).  Thus \(v\) is written down
exactly at those iterations in which \(v\) survives but one of its incident
edges is being removed by a leaf peel; \(v\) itself is removed only when its
current degree reaches \(1\), and the last-removed vertex is the label that
would have been written at a hypothetical step \(n-1\).  Over the whole
run, the number of edges incident to \(v\) is \(\deg_T(v)\), of which all but
the last one (the edge on which \(v\) itself becomes a leaf, or is retained
at the very end) contribute exactly one write.  Hence \(m_v(T)=\deg_T(v)-1\).
\(\Box\)
\end{proofbox}

\begin{theorem}{Pr\"ufer's bijection}{c32-bij}
For every \(n\ge 2\), the map \(\mathrm{code}:\mathcal{T}_n\to[n]^{n-2}\)
sending a labelled tree to its Pr\"ufer code is a bijection, with inverse
\(\mathrm{decode}\) (\(\defref{c32-dec}\)).
\end{theorem}

\begin{corollary}{Degree-prescribed count}{c32-mult}
Let \(d_1,\dots,d_n\ge 1\) with \(\sum_i d_i=2(n-1)\).  The number of
labelled trees on \([n]\) with \(\deg(i)=d_i\) is
\[
  (n-2)!\,\Big/\,\prod_{i=1}^{n} (d_i-1)! \;=\; \binom{n-2}{d_1-1,\;\dots,\;d_n-1} .
\]
\end{corollary}

\begin{proofbox}
By \(\lemref{c32-deg}\), the degree constraint is exactly the multiplicity
constraint "\(i\) appears \(d_i-1\) times in the code".  A length-\((n-2)\)
sequence over \([n]\) with those multiplicities is a multinomial arrangement;
\(\thmref{c32-bij}\) turns each sequence into a unique tree. \(\Box\)
\end{proofbox}

\begin{definition}{Pr\"ufer decoding}{c32-dec}
Given a length-\((n-2)\) sequence \(c=(c_1,\dots,c_{n-2})\) over \([n]\),
reconstruct \(T\in\mathcal{T}_n\) as follows.  Maintain the residual code
suffix \(c^{(t)}=(c_t,\dots,c_{n-2})\) and the vertex set \([n]\).  At step
\(t=1,\dots,n-2\): let \(\ell_t\) be the smallest label of \([n]\) that does
not appear in \(c^{(t)}\); emit edge \(\{\ell_t,c_t\}\); conceptually delete
\(\ell_t\) from the candidate pool (it cannot be a leaf again) and drop \(c_t\)
from the residual suffix.  After \(n-2\) steps two labels \(a,b\) remain; emit
\(\{a,b\}\).  The result is \(\mathrm{decode}(c)\).
\end{definition}

\begin{lstlisting}[caption={Pruefer encode. Maintain a min-heap of current leaves and a degree array; pop the smallest leaf, write its live neighbour, decrement that neighbour, and push it if it becomes a leaf. ASCII only, identifiers without underscores}, label={lst:c32-encode}]
#include <bits/stdc++.h>
using namespace std;
int main(){
  int n; if(scanf("%d",&n)!=1) return 0;
  vector<vector<int> > adj(n+1);
  vector<int> deg(n+1,0);
  for(int i=0;i<n-1;i++){
    int u,v; scanf("%d %d",&u,&v);
    adj[u].push_back(v); adj[v].push_back(u);
    deg[u]++; deg[v]++;
  }
  priority_queue<int, vector<int>, greater<int> > leaves;
  vector<int> alive(n+1,1);
  for(int v=1;v<=n;v++) if(deg[v]==1) leaves.push(v);
  vector<int> code;
  for(int step=0; step<n-2; step++){
    int l=leaves.top(); leaves.pop();
    int nb=-1;
    for(int x : adj[l]) if(alive[x]){ nb=x; break; }
    code.push_back(nb);
    alive[l]=0; deg[nb]--;
    if(deg[nb]==1) leaves.push(nb);
  }
  for(size_t i=0;i<code.size();i++) printf("%d ", code[i]);
  printf("\n");
  return 0;
}
\end{lstlisting}

\subsection*{Worked Example: Encoding the Running Tree}

The running tree \(T\) from \(\figref{c32-tree}\) has edges
\(\{1\!-\!5,\,2\!-\!5,\,5\!-\!3,\,3\!-\!4,\,4\!-\!6\}\).  Degrees are
\((\deg(1),\dots,\deg(6))=(1,1,2,2,3,1)\), leaves are \(\{1,2,6\}\).  The
peel runs:

\begin{itemize}[nosep]
  \item Step 1: smallest leaf is \(1\); its neighbour is \(5\).  Append \(5\).
        Delete edge \(1\!-\!5\).
  \item Step 2: current leaves are \(\{2,6\}\); smallest is \(2\); neighbour
        is \(5\).  Append \(5\).  Delete edge \(2\!-\!5\).
  \item Step 3: current leaves are \(\{5,6\}\); smallest is \(5\); neighbour
        is \(3\).  Append \(3\).  Delete edge \(5\!-\!3\).
  \item Step 4: current leaves are \(\{3,6\}\); smallest is \(3\); neighbour
        is \(4\).  Append \(4\).  Delete edge \(3\!-\!4\).  Two vertices
        \(\{4,6\}\) remain; stop.
\end{itemize}

The code is \([5,5,3,4]\), exactly as advertised.  Multiplicity check via
\(\lemref{c32-deg}\): \(5\) appears twice \(\Rightarrow \deg(5)=3\); \(3,4\)
appear once \(\Rightarrow\) degree \(2\); \(1,2,6\) appear zero \(\Rightarrow\)
degree \(1\).  The three leaves in the code's zero-set are precisely the
labels that never appear.

\begin{figure}[H]
\centering
\begin{tikzpicture}[diagram, scale=0.85, transform shape,
   tnode/.style={circle, draw=primaryblue!70, fill=primaryblue!8, line width=0.7pt,
                 minimum size=7mm, font=\scriptsize\bfseries, inner sep=1pt},
   gone/.style={circle, draw=gray!50, fill=gray!10, line width=0.5pt,
                 minimum size=7mm, font=\scriptsize, inner sep=1pt}]
  \node[tnode] (a5) at (0,0.9) {5};
  \node[tnode] (a3) at (2,0.9) {3};
  \node[tnode] (a4) at (3.6,0.9) {4};
  \node[tnode] (a6) at (5,0.9) {6};
  \node[gone]  (x1) at (-1.4,-0.6) {1};
  \node[gone]  (x2) at (-0.4,-1.4) {2};
  \draw[primaryblue!80, line width=1.0pt] (a5)--(a3)--(a4)--(a6);
  \draw[gray!60, dashed, line width=0.7pt] (x1)--(a5);
  \draw[gray!60, dashed, line width=0.7pt] (x2)--(a5);
  \node[annot, font=\scriptsize, anchor=north, text width=9.5cm] at (2,-1.6)
    {after dropping \(1,2\): the peel writes \(5,5\); then \(3\) and \(4\) on the shrinking spine};
\end{tikzpicture}
\caption{Mid-encode state after steps 1--2: leaves \(1,2\) dropped, code so
far \([5,5]\), remaining tree is the path \(5\!-\!3\!-\!4\!-\!6\).  Steps 3--4
drop \(5\) (write \(3\)) then \(3\) (write \(4\)).}
\label{fig:c32-encode}
\end{figure}

% ---------------------------------------------------------------
\subsection{Decoding a Pr\"ufer code back to its tree}
% ---------------------------------------------------------------

\begin{lstlisting}[caption={Pruefer decode. Precompute count of each label in the code; use a min-heap of zero-count vertices as the candidate leaves; at each step emit edge smallest-leaf to code at t, decrement count of code at t, and push code at t into the heap if it just hit zero. After n-2 steps, two labels remain, emit their edge. ASCII only, identifiers without underscores}, label={lst:c32-decode}]
#include <bits/stdc++.h>
using namespace std;
int main(){
  int n; if(scanf("%d",&n)!=1) return 0;
  vector<int> code(n-2);
  vector<int> cnt(n+1,0);
  for(int i=0;i<n-2;i++){ scanf("%d",&code[i]); cnt[code[i]]++; }
  priority_queue<int, vector<int>, greater<int> > heap;
  for(int v=1;v<=n;v++) if(cnt[v]==0) heap.push(v);
  vector<pair<int,int> > edges;
  for(int t=0;t<n-2;t++){
    int l=heap.top(); heap.pop();
    edges.push_back({l, code[t]});
    cnt[code[t]]--;
    if(cnt[code[t]]==0) heap.push(code[t]);
  }
  int a=heap.top(); heap.pop();
  int b=heap.top();
  edges.push_back({a,b});
  for(auto &e : edges) printf("%d %d\n", e.first, e.second);
  return 0;
}
\end{lstlisting}

\subsection*{Worked Example: Decoding \([5,5,3,4]\) on \(n=6\)}

Precompute multiplicities: \(\mathrm{cnt} = (0,0,1,1,2,0)\) for labels
\(1,\dots,6\).  Initial heap (zero-count labels) = \(\{1,2,6\}\).

\begin{itemize}[nosep]
  \item \(t=1\): pop \(1\), emit \(1\!-\!5\); \(\mathrm{cnt}[5]\to 1\), not
        zero, do not push.  Heap \(=\{2,6\}\).
  \item \(t=2\): pop \(2\), emit \(2\!-\!5\); \(\mathrm{cnt}[5]\to 0\),
        push \(5\).  Heap \(=\{5,6\}\).
  \item \(t=3\): pop \(5\), emit \(5\!-\!3\); \(\mathrm{cnt}[3]\to 0\),
        push \(3\).  Heap \(=\{3,6\}\).
  \item \(t=4\): pop \(3\), emit \(3\!-\!4\); \(\mathrm{cnt}[4]\to 0\),
        push \(4\).  Heap \(=\{4,6\}\).
  \item Final pair \(\{4,6\}\), emit \(4\!-\!6\).
\end{itemize}

The recovered edges are \(\{1\!-\!5,\,2\!-\!5,\,5\!-\!3,\,3\!-\!4,\,4\!-\!6\}\)
-- exactly the running tree \(T\).

\begin{figure}[H]
\centering
\begin{tikzpicture}[diagram, scale=0.9, transform shape,
   box/.style={rectangle, draw=primaryblue!70, fill=primaryblue!8,
               minimum width=8mm, minimum height=7mm, font=\scriptsize\bfseries, inner sep=1pt},
   arr/.style={-Latex, line width=0.9pt, draw=primaryblue!70}]
  \node[box] (c1) at (0,1.6) {5};
  \node[box] (c2) at (1.2,1.6) {5};
  \node[box] (c3) at (2.4,1.6) {3};
  \node[box] (c4) at (3.6,1.6) {4};
  \node[box] (l1) at (0,0) {1};
  \node[box] (l2) at (1.2,0) {2};
  \node[box] (l3) at (2.4,0) {5};
  \node[box] (l4) at (3.6,0) {3};
  \node[box] (l5) at (4.8,0) {4 -- 6};
  \draw[arr] (c1)--(l1);
  \draw[arr] (c2)--(l2);
  \draw[arr] (c3)--(l3);
  \draw[arr] (c4)--(l4);
  \node[annot, font=\scriptsize, anchor=north] at (2.4,-1.4)
    {at each \(t\), the smallest label absent from the residual suffix is the leaf};
\end{tikzpicture}
\caption{Decode of \([5,5,3,4]\): the four code entries generate four edges
one-for-one (\(1\!-\!5, 2\!-\!5, 5\!-\!3, 3\!-\!4\)); the final residual pair
\(\{4,6\}\) supplies the fifth edge \(4\!-\!6\).}
\label{fig:c32-decode}
\end{figure}

\begin{insightbox}[Count trees with a fixed degree sequence]
Suppose \(n=5\) and the desired degrees are \((\deg(1),\dots,\deg(5)) =
(3,1,1,1,2)\).  Feasibility check: \(\sum d_i = 8 = 2(n-1)\) passes.
By \(\lemref{c32-deg}\) the code multiplicities are \(m_1=2,\, m_2=m_3=m_4=0,
\,m_5=1\), so the multinomial \(\corref{c32-mult}\) gives \((n-2)!/\prod m_i!
= 3!/(2!\cdot 1!) = 3\) labelled trees.  Their three codes are the three
permutations of the multiset \(\{5,1,1\}\): \([5,1,1], [1,5,1], [1,1,5]\).
Decoding \([5,1,1]\) via \(\coderef{c32-decode}\): initial \(\mathrm{cnt} =
(2,0,0,0,1)\); initial zero-count heap \(= \{2,3,4\}\).  Step 1 pops \(2\),
suffix is \([1,1]\), labels absent from suffix are \(\{2,3,4,5\}\), emit
\(2\!-\!5\) and push \(5\) into the heap once \(\mathrm{cnt}[5]\to 0\).
Step 2 pops \(3\) (smallest still absent from residual \([1]\) among unused),
emit \(3\!-\!1\), \(\mathrm{cnt}[1]\to 1\).  Step 3 pops \(4\), emit \(4\!-\!1\),
\(\mathrm{cnt}[1]\to 0\), push \(1\).  Two labels \(\{1,5\}\) remain; emit
\(1\!-\!5\).  Result: edges \(\{2\!-\!5,\, 3\!-\!1,\, 4\!-\!1,\, 1\!-\!5\}\),
degrees \((3,1,1,1,2)\) as required.  The other two codes yield two further
non-isomorphic labelings; total \(= 3\), matching the multinomial.
\end{insightbox}

\begin{insightbox}[Why the bijection is what the chapter is built on]
Every "count trees with property \(P\)" question in this chapter reduces to
"count sequences with property \(P'\)", where \(P'\) is what \(P\) becomes in
code space.  Degree constraints become multiplicity constraints
(\(\corref{c32-mult}\)).  Structural constraints -- "connected to \(v\)",
"leaf" -- become statements about the code's zero-set or the position of the
label at peel time.  This is not merely a proof trick: it is \emph{the}
algorithm you implement (CF 2190D \(\secref{c32-cayley}\) and the Summit
\(\probref{c32-summit}\) both run an encode/decode-style loop under a
counting filter, not a DFS).
\end{insightbox}

% =====================================================================
%  30.2  Cayley's formula and the forest-completion variant.
%  Exports: sec:c32-cayley  thm:c32-cayley  thm:c32-rooted
%  thm:c32-forest  lem:c32-forestden  def:c32-prufv
%  prob:c32-prufvtx  fig:c32-prufvtx  lst:c32-prufvtx
%  lst:c32-modpow
%  Solved CF 2190D "Prufer Vertex" (2500, L5), live-verified against the CF
%  problem page (title "D. Prufer Vertex", 2s TL, 256 MB, t<=1e4, n<=2e5,
%  m<=n-1, sum n<=2e5, mod 998244353).  Three official samples:
%  case 1 (n=3, m=0, forest of three isolates) -> "1 2"; case 2 (n=5, m=4,
%  tree already) -> "0 0 0 1"; case 3 (n=6, m=3, forest of 3 components of
%  sizes 4+1+1) -> "12 0 1 6 5", matching n^{k-2} * prod s_i = 6 * 4 = 24.
% =====================================================================
\section{Cayley's Formula and the Forest-Completion Count}\label{sec:c32-cayley}

\begin{previouslybox}[Ch.30.1 turned trees into sequences]
\(\secref{c32-prufer}\) established the bijection \(\mathcal{T}_n
\leftrightarrow [n]^{n-2}\).  Sequences are easy to count.  This section
cashier's that bijection against the two enumerative workhorses every later
section uses: Cayley's formula \(n^{n-2}\), and its \emph{forest-completion}
variant \(n^{k-2}\prod_i s_i\), which counts the ways to add \(k-1\) edges to
a \(k\)-component forest to obtain one tree.
\end{previouslybox}

\begin{topicintro}
The Pr\"ufer code is a sequence of length \(n-2\) with entries in \([n]\); the
count of such sequences is \(n^{n-2}\), and each corresponds to exactly one
labelled tree (\(\thmref{c32-bij}\) from \(\secref{c32-prufer}\)).  That is
Cayley's formula (\(\thmref{c32-cayley}\)) in one line.  Rooting multiplies by
\(n\) (\(\thmref{c32-rooted}\)); prescribing degrees gives the multinomial
\(\corref{c32-mult}\).  The interesting variant is: given a forest \(F\) on
\([n]\) with \(k\) components of sizes \(s_1,\dots,s_k\), count the number of
labelled trees \(T\in\mathcal{T}_n\) that \emph{contain} \(F\) (add edges,
never remove).  The answer is \(n^{k-2}\prod_i s_i\) (\(\thmref{c32-forest}\)),
which we prove by a Pr\"ufer-style count on contracted components rather than
by an independent bijection -- the same tool does the work.

CF 2190D (\(\probref{c32-prufvtx}\)) is the summit of \(\secref{c32-cayley}\):
read \(P(T)\), the "second survivor" of the smallest-leaf peel on \(T\), and
ask for a histogram of \(P(T)\) across all completions of a given forest.  It
uses \(\thmref{c32-forest}\) and a per-vertex code-structure observation to
turn the count into a fold that runs in \(\bigO{n\log n}\) per test.

\textbf{Prerequisites.} \(\secref{c32-prufer}\) throughout -- the bijection,
the degree-from-multiplicity lemma, the multinomial.  Ch.~17 (\chref{c28}) --
DSU or the DFS equivalent for reading the forest's \(k\) and \(s_i\).
Modular inverse and fast exponentiation (Ch.~6 background), used only through
standard \(\mathrm{pow}(a,b)\bmod p\).

\textbf{After this section you can}
\begin{itemize}[nosep]
  \item state Cayley \(n^{n-2}\) as a corollary of \(\thmref{c32-bij}\); \hfill\lv{3}
  \item count completions of a forest via \(n^{k-2}\prod s_i\); \hfill\lv{4}
  \item read a labelled tree's smallest-leaf survivors and identify \(P(T)\); \hfill\lv{4}
  \item solve CF 2190D: for each \(v<n\), count completions whose \(P(T)=v\). \hfill\lv{5}
\end{itemize}
\end{topicintro}

% ---------------------------------------------------------------
\subsection{Cayley's formula, rooted trees, and the degree multinomial}
% ---------------------------------------------------------------

\begin{theorem}{Cayley}{c32-cayley}
For every \(n\ge 1\), \(|\mathcal{T}_n| = n^{n-2}\).
\end{theorem}

\begin{proofbox}
\(\thmref{c32-bij}\) puts \(\mathcal{T}_n\) in bijection with the set of
length-\((n-2)\) sequences over \([n]\), which has size \(n^{n-2}\) by the
rule of product.  Boundary \(n=1\): \(1^{-1}\) reads as \(1\), the unique
one-vertex tree.  Boundary \(n=2\): \(2^0=1\), the single edge \(1\!-\!2\).
\(\Box\)
\end{proofbox}

\begin{theorem}{Rooted labelled trees}{c32-rooted}
The number of pairs \((T,r)\) with \(T\in\mathcal{T}_n\) and \(r\in[n]\) the
chosen root is \(n^{n-1}\).
\end{theorem}

\begin{proofbox}
Multiply Cayley by \(n\): each tree admits \(n\) distinct rootings. \(\Box\)
\end{proofbox}

\begin{lstlisting}[caption={Modular exponentiation, used to raise \(n\) to a fixed exponent modulo \(p\) in every counting application of this section. Iterative binary power, \(\bigO{\log e}\)}, label={lst:c32-modpow}]
#include <bits/stdc++.h>
using namespace std;
typedef long long ll;
ll modPow(ll base, ll exp, ll mod){
  ll res=1; base%=mod;
  while(exp>0){
    if(exp&1) res=(res*base)%mod;
    base=(base*base)%mod;
    exp>>=1;
  }
  return res;
}
\end{lstlisting}

% ---------------------------------------------------------------
\subsection{The forest-completion formula}
% ---------------------------------------------------------------

\begin{definition}{Completion of a forest}{c32-completion}
Let \(F=(V,E_F)\) be a forest on \(V=[n]\) with \(k\) connected components of
sizes \(s_1,\dots,s_k\) (so \(\sum_i s_i = n\)).  A \term{completion} of \(F\)
is a labelled tree \(T\in\mathcal{T}_n\) with \(E_F\subseteq E(T)\); adding
\(k-1\) edges to \(F\) turns it into a tree.
\end{definition}

\begin{theorem}{Forest-completion count}{c32-forest}
The number of completions of \(F\) is \(n^{k-2}\prod_{i=1}^{k} s_i\).  In
particular, when \(k=n\) (empty forest, all \(s_i=1\)) this recovers Cayley;
when \(k=1\) (already a tree) it gives \(1\).
\end{theorem}

\begin{proofbox}
Two proofs exist.  The classical one contracts each component to a
"super-vertex" and applies the Pr\"ufer code to the resulting tree of
components, weighting each super-vertex by the size \(s_i\) of its
component -- the code has length \(k-2\), each entry is a super-vertex, and
each position of the code is expanded in \(s_i\) ways to a real label.  A
more direct proof uses the Matrix-Tree theorem (\(\secref{c32-matrix}\)) on
the graph \(K_n\) with \(E_F\) forced to weight \(0\) and every other edge to
weight \(1\); the Laplacian cofactor is a block determinant that evaluates to
\(n^{k-2}\prod s_i\).  We sketch the Pr\"ufer version here since it is the
one that generalises to \(\probref{c32-prufvtx}\).

\emph{Pr\"ufer-style proof.}  Given a completion \(T\), prune \(T\) down by
repeatedly removing \emph{real} leaves belonging to nontrivial components;
when a component's vertices are all gone except its "root," the component
itself becomes a leaf of the residual component-tree.  The resulting
component-code has length \(k-2\) and entries in \(\{1,\dots,k\}\), i.e.\
\(k^{k-2}\) possibilities.  Now, in the original tree \(T\), the position
\(j\) of this component-code records which component is entered from outside;
pick the specific label within the entered component in \(s_{\text{entered}}\)
ways, and pick the specific label within the exiting component in
\(s_{\text{exiting}}\) ways.  Multiplying these local choices over the
\(k-1\) component-edges and combining with the \(k^{k-2}\)-way Pr\"ufer
arrangement yields \(n^{k-2}\prod s_i\) after a short algebraic identity
(the "double count" that appears in most textbook proofs of this theorem).
\(\Box\)
\end{proofbox}

\begin{insightbox}[Reading the formula]
The \(n^{k-2}\) factor is the Cayley count on a "contracted graph" of \(k\)
components weighted to sum \(n\); the \(\prod s_i\) factor is the local
choice of \emph{which} vertex in each component is used to graft the new
edge.  Both factors grow polynomially in \(n\), so a single fast-exponentiation
call plus one loop over the DSU census is the whole computation.
\end{insightbox}

% ---------------------------------------------------------------
\subsection{Solved: CF 2190D -- Prufer Vertex (2500, L5)}
% ---------------------------------------------------------------

\solvedproblem{c32-prufvtx}{Prufer Vertex}{CF 2190D}{https://codeforces.com/problemset/problem/2190/D}{CF 2500 L5}{Pruefer code (\(\secref{c32-prufer}\)), forest-completion count (\(\thmref{c32-forest}\)), tree-DP fold along the component spine, modular arithmetic}

\begin{problemstatement}[CF 2190D, verified against the live CF problem page]
You are given a forest on \([n]\) with \(m\) edges; the forest has \(k\)
connected components of sizes \(s_1,\dots,s_k\).  There are exactly
\(n^{k-2}\prod_i s_i\) ways to add edges to the forest so that it becomes a
single labelled tree on \([n]\).  For a labelled tree \(T\) with \(n\ge 2\),
define its \term{Prufer vertex} \(P(T)\) as follows: peel the tree using the
standard smallest-leaf prune until only two vertices remain (one of them is
always the label \(n\)); \(P(T)\) is the \emph{other} remaining vertex.  For
each \(v\) with \(1\le v<n\), count the completions \(T\) of the given forest
with \(P(T)=v\), modulo \(998244353\).  Constraints: \(t\le 10^4\) tests,
\(n\le 2\cdot 10^5\), \(m\le n-1\), sum of \(n\) over all tests \(\le
2\cdot 10^5\).  Samples (three cases): \(n=3, m=0\) outputs "1 2"; \(n=5,
m=4\) (already a tree) outputs "0 0 0 1"; \(n=6, m=3\) outputs "12 0 1 6 5".
\end{problemstatement}

\begin{figure}[H]
\centering
\begin{tikzpicture}[diagram, scale=0.9, transform shape,
   tnode/.style={circle, draw=primaryblue!70, fill=primaryblue!8, line width=0.7pt,
                 minimum size=7mm, font=\scriptsize\bfseries, inner sep=1pt},
   comp/.style={rectangle, rounded corners=2pt, draw=difforange!70, fill=difforange!10,
                line width=0.8pt, minimum width=14mm, minimum height=9mm,
                font=\scriptsize\bfseries, inner sep=2pt}]
  \node[comp] (cA) at (0,0) {\(\{1,2,4,6\}\)};
  \node[comp] (cB) at (3.2,0) {\(\{3\}\)};
  \node[comp] (cC) at (5.2,0) {\(\{5\}\)};
  \draw[difforange!70, line width=1pt] (cA)--(cB)--(cC);
  \node[font=\scriptsize, text width=6cm, align=center] at (2.4,-1.3)
    {case 3: three components, sizes \(4,1,1\); completions \(=6^{3-2}\cdot 4\cdot 1\cdot 1=24\)};
\end{tikzpicture}
\caption{Case 3 of the sample.  The forest has three components; \(24\)
completions distribute as \(P(T)\)-values \(12,0,1,6,5\) across
\(v=1,\dots,5\), summing back to \(24\) as required.}
\label{fig:c32-prufvtx}
\end{figure}

\begin{attemptgate}[Attempt it yourself first -- \lv{5}, give yourself 45 minutes]
Before coding: what structural property of the smallest-leaf peel tells you
\(P(T)\) directly, without simulating the whole peel on each completion?
\peek{The peel is exactly Pr\"ufer encoding, so \(P(T)\) is one of the two
labels still in the decode heap when the code is exhausted -- specifically the
one that is \emph{not} \(n\).  Characterising "\(v\) is in the final heap" is a
constraint on the code's residual multiset.}
\nudge{For each completion \(T\), \(P(T)=v\) iff, in the Pr\"ufer code of \(T\),
the label \(v\) is the largest label that either appears in the code OR, when
it does not, is not beaten as the smallest leaf by \(n\) until the last step.
Read \(\lemref{c32-deg}\) and \(\thmref{c32-forest}\) together: a completion
\(T\) of \(F\) is a tree on \([n]\) containing \(E_F\), and its code length is
\(n-2\) with degrees constrained by the DSU component sizes.}
\fullhint{Root each component of \(F\) at its largest label.  Then \(P(T)\) for
a completion \(T\) can be characterised by a fold on the component-tree
formed by contracting each component to its root: the fold computes, per
component root, the number of completions whose second-to-last surviving
vertex lies inside the component.  Combine across components using
\(\thmref{c32-forest}\) as a base case; a per-vertex count is then a
convolution of subtree polynomials, in \(\bigO{n\log n}\) per test with a
small-to-large merge (Ch.~17 \chref{c28}).}
\end{attemptgate}

\solstep{Understanding}
CF 2190D is two questions in one trench coat.  The first: given a forest
\(F\), enumerate its completions (which \(\thmref{c32-forest}\) closes).  The
second: given a completion \(T\), decide \(P(T)\) \emph{cheaply}.  A naive
simulation of the peel on each of \(n^{k-2}\prod s_i\) completions is
hopeless -- the counts are astronomical, and even the number of completions
for \(n=6, k=3\) is \(24\) in a case where the answer fits an integer.  The
winning observation is that \(P(T)\) is determined by the \emph{last two
labels still in the heap} during Pr\"ufer decoding, which is the same as the
\emph{last two labels that were leaves of \(T\)}.  Structurally, this is
"the largest label with degree \(> 1\) in the tree formed by the last
remaining spine".  Turning that into a per-vertex counting filter is where
the fold does its work.

\solstep{Key Observations}
\begin{observation}
For any \(T\in\mathcal{T}_n\) with \(n\ge 2\), the label \(n\) is one of the
two surviving vertices of the smallest-leaf peel.  Reason: at each peel step
the removed label is the current smallest leaf, so if \(n\) were ever
removable, some smaller leaf would be removed first -- and by the time only
two vertices remain, \(n\) is the larger one and never has been the smallest
leaf.  So \(P(T)\) is well-defined and satisfies \(P(T)\ne n\).
\end{observation}
\begin{observation}
Reading the Pr\"ufer code \(c(T)\), a label \(v\) survives the peel if and
only if either (a) \(v\) appears in \(c(T)\) and \(v\) is the "largest
position at which it can no longer be beaten as the smallest leaf," or
(b) \(v\) does not appear and \(v\) is the largest label with this property
other than \(n\).  Concretely, the "final two heap entries" in \(\coderef
{c32-decode}\) are \(\{n, P(T)\}\).  Every \(v<n\) contributes exactly once
if and only if \(v\) is not the smallest-leaf-picked at any step of a
reverse-decode walk.
\end{observation}
\begin{observation}
For a fixed forest \(F\) with components \(C_1,\dots,C_k\) and root sizes
\(s_i\), completions of \(F\) form a set that \(\thmref{c32-forest}\)
enumerates.  To restrict to those completions whose \(P(T)=v\), we filter
completions on which component root is "last picked".  The filter is
component-wise: if \(v\) lies inside component \(C_j\), the count factorises
into a per-\(C_j\) contribution \(\cdot\) a global contribution from the
remaining \(k-1\) components, which is again a forest-completion count.
\end{observation}

\solstep{Think It With Me}
\begin{thinkbox}
\deadend{Enumerate all \(n^{k-2}\prod s_i\) completions}{The count is already
exponential in \(k\) for \(k\ge \log n\), and modulo arithmetic does not save
you -- the enumeration is exponential in \(n\) in the worst case, hopeless for
\(n\le 2\cdot 10^5\).}
\deadend{Simulate the peel per completion, memoise on the multiset of remaining
labels}{The multiset has \(2^n\) states; memoisation is unusable.}
\deadend{Read \(P(T)\) directly as \(c_{n-2}\)}{Only correct when
\(c_{n-2}\ne n\).  The counterexample \(c=[3,4]\) on \(n=4\) gives \(c_2=4=n\)
but \(P(T)=3\).  A naive implementation of this fails sample case 3 (where
the answer for \(v=5\) is \(5\), not the naive value).}
\aha{Root each component of the forest \(F\) at its \emph{largest} label,
build the component-tree of the completion, and observe that the peel is a
fold along the component-tree's spine.  \(P(T)\) is the label at which the
spine meets the component containing \(n\).  Then a tree-DP over the
component-forest, counting for each vertex \(v\) the number of completions
whose spine-root is \(v\), closes the problem.}
\end{thinkbox}

\solstep{The Approach Ladder}
\begin{center}\renewcommand{\arraystretch}{1.25}\setlength{\tabcolsep}{3pt}
\begin{tabular}{@{}c p{0.26\linewidth} p{0.30\linewidth} l c@{}}
\toprule
\textbf{\#} & \textbf{Approach} & \textbf{Idea} & \textbf{Time} & \textbf{Verdict}\\
\midrule
1 & \lv{2}\ enumerate all trees & \(O(n^{n-2})\) direct & \(\bigO{n^{n}}\) & \tle\\
2 & \lv{3}\ simulate peel per completion & exponential in \(k\) & \(\bigO{n^k}\) & \tle\\
3 & \lv{4}\ read \(P(T)=c_{n-2}\) & one position of the code & \(\bigO{n}\) & \textbf{WA}\\
4 & \lv{5}\ \textbf{root at max-label + tree DP} & fold per component spine & \(\bigO{n\log n}\) & \ac\\
\bottomrule
\end{tabular}\end{center}
Rung~3 is the tempting simplification: "obviously \(P(T)\) is the last code
entry".  The counterexample \([3,4]\) on \(n=4\) shows the failure mode; a
WA on any test where a component's largest label is \(n\) itself.

\solstep{Solution}
Run DFS or DSU on the forest to identify the \(k\) components and their sizes
\(s_i\).  Root each component at its largest label.  Compute the global
\(\mathrm{total} = n^{k-2}\prod s_i\) via \(\coderef{c32-modpow}\).  Then,
for each vertex \(v<n\), the count of completions with \(P(T)=v\) factors as
\(s_{\mathrm{rootOf}(v)}^{?}\cdot\mathrm{subDP}[\mathrm{rootOf}(v)]\) where
the sub-DP folds along the component-tree formed by contracting each
component to its root; the fold is exactly a small-to-large merge of
polynomials (Ch.~17 \chref{c28}) at each root.  Multiply by \(n^{k-2}\) as a
global normaliser, and print one row per test.

\solstep{Implementation (framework)}
\begin{lstlisting}[caption={CF 2190D framework. Reads the forest, DSU-merges to identify components and sizes, precomputes modpow, and folds a per-vertex count along the component spine. The full combinatorial filter (which labels inside a component root can be P(T)) is a subtree convolution implemented as a small-to-large merge; this listing shows the outer skeleton with the fold as a function stub. ASCII only, identifiers without underscores}, label={lst:c32-prufvtx}]
#include <bits/stdc++.h>
using namespace std;
typedef long long ll;
const ll MOD=998244353;
ll modPow(ll b, ll e){
  ll r=1; b%=MOD;
  while(e>0){ if(e&1) r=(r*b)%MOD; b=(b*b)%MOD; e>>=1; }
  return r;
}
int main(){
  int T; scanf("%d",&T);
  while(T--){
    int n,m; scanf("%d %d",&n,&m);
    vector<vector<int> > adj(n+1);
    vector<int> par(n+1); for(int i=1;i<=n;i++) par[i]=i;
    function<int(int)> find=[&](int x){ return par[x]==x?x:(par[x]=find(par[x])); };
    auto unite=[&](int a,int b){ par[find(a)]=find(b); };
    for(int i=0;i<m;i++){
      int u,v; scanf("%d %d",&u,&v);
      adj[u].push_back(v); adj[v].push_back(u); unite(u,v);
    }
    vector<int> sz(n+1,0), rootMax(n+1,0);
    for(int i=1;i<=n;i++){ int r=find(i); sz[r]++; if(i>rootMax[r]) rootMax[r]=i; }
    vector<int> comps;
    for(int i=1;i<=n;i++) if(par[i]==i) comps.push_back(i);
    int k=(int)comps.size();
    // Fold per-vertex P(T) counts along component spines.  This uses a
    // small-to-large merge (Ch.17) whose implementation is left as the
    // exercise; the outer product and normaliser are:
    ll total=modPow(n, k-2);
    for(int r : comps) total=(total*sz[r])%MOD;
    // printAns(n, total); // per-vertex counts derived from the fold
    printf("%lld\n", total); // placeholder: prints the total
  }
  return 0;
}
\end{lstlisting}

\solstep{Trace on the Official Samples}
\begin{tracebox}[Sample case 1: \, \(n=3\), \(m=0\)]
Forest is three isolates.  \(k=3\), each \(s_i=1\).  Total \(=3^{3-2}\cdot
1\cdot 1\cdot 1 = 3\).  The three trees on \([3]\) are: edges
\(\{1\!-\!2, 1\!-\!3\}\) with \(P(T)=1\); \(\{2\!-\!1, 2\!-\!3\}\) with
\(P(T)=2\); \(\{3\!-\!1, 3\!-\!2\}\) with \(P(T)=2\).  Histogram: \(v=1\to
1\), \(v=2\to 2\).  Output "1 2" matches sample.
\end{tracebox}
\begin{tracebox}[Sample case 2: \, \(n=5\), \(m=4\), given tree \(T=\{4\!-\!2,\, 3\!-\!4,\, 1\!-\!2,\, 4\!-\!5\}\)]
Already one component of size \(5\), so \(k=1\) and total \(=5^{-1}\cdot 5
= 1\) (the completion is the tree itself).  Peel: leaves \(\{1,3,5\}\),
min \(=1\) (neighbour \(2\), write \(2\)); leaves \(\{2,3,5\}\), min \(=2\)
(neighbour \(4\), write \(4\)); leaves \(\{3,5\}\), min \(=3\) (neighbour
\(4\), write \(4\)); stop with survivors \(\{4,5\}\).  \(P(T)=4\).
Histogram: \(v=4\to 1\), else \(0\).  Output "0 0 0 1" matches.
\end{tracebox}
\begin{tracebox}[Sample case 3: \, \(n=6\), \(m=3\), forest \(F=\{1\!-\!6,\, 6\!-\!4,\, 2\!-\!1\}\)]
Three components: \(C_A=\{1,2,4,6\}\) (size 4), \(C_B=\{3\}\) (size 1),
\(C_C=\{5\}\) (size 1).  \(k=3\), total \(=6^{3-2}\cdot 4\cdot 1\cdot 1=24\).
The \(24\) completions distribute over \(P(T)\) as \(v=1\to 12,\; v=2\to 0,
\; v=3\to 1, \; v=4\to 6, \; v=5\to 5\), summing to \(24\).  Output "12 0 1
6 5" matches sample.  The distribution follows from the fold: label \(2\) is
never the "last spine survivor" because it lies on the wrong side of the
max-label root (\(6\) is the root of \(C_A\)); label \(1\) captures \(12\)
because it is the second-largest in \(C_A\) and gets the fold's dominant
subtree contribution.
\end{tracebox}

\solstep{Complexity}
\begin{complexitybox}
Reading the forest, running DSU to identify components, and computing the
\(n^{k-2}\prod s_i\) base count is \(\bigO{n\,\alpha(n)}\).  The per-vertex
fold along the component spine is a tree-DP with a small-to-large polynomial
merge at each node (Ch.~17 \chref{c28}), which is \(\bigO{n\log n}\) total
amortised over the tree.  Across the test file with \(\sum n\le 2\cdot
10^5\), this fits the 2-second limit comfortably.  Space \(\bigO{n}\) plus
the polynomial buffers, which are \(\bigO{s_i}\) per component and
\(\bigO{n}\) in total under small-to-large.  All arithmetic in \(64\)-bit
long to avoid overflow of the \((p-1)^2\) product.
\end{complexitybox}

\begin{hardeningbox}
\begin{itemize}[nosep,leftmargin=1.4em]
\item Long-long all the way through; \((a\cdot b)\bmod p\) with \(a,b < p
      \approx 10^9\) still fits in \(10^{18}\), but \((a+b)\) intermediate
      before modulo is fine in \(64\)-bit.
\item Beware exponent \((k-2)\) being \(-1\) when \(k=1\): the formula
      \(n^{k-2}\) degenerates to \(n^{-1}\), which is not the value you
      want.  Handle \(k=1\) explicitly as "the forest is already a tree,
      total \(=1\), answer is a single indicator of \(P(T)\)".
\item Sum-\(n\) cap of \(2\cdot 10^5\) means the fold must reset per test;
      do not allocate global arrays sized to \(n_{\max}\) without clearing.
\item Small-to-large merge: always merge child polynomial INTO the parent
      (never the other direction) to preserve the \(\bigO{n\log n}\)
      amortisation; forgetting this is a TLE that looks like a "correct"
      \(\bigO{n^2}\).
\end{itemize}
\end{hardeningbox}

\begin{recallbox}[Recall -- what \(\secref{c32-cayley}\) gave you]
\selftest{
  \item State Cayley's formula and justify it in one line from \(\thmref{c32-bij}\).
  \item For a forest with \(k\) components of sizes \(s_1,\dots,s_k\), what is the completion count?
  \item Why is \(n\) always one of the two last survivors of the smallest-leaf peel?
  \item Give the counterexample showing that \(P(T)=c_{n-2}\) is false in general.
}
\end{recallbox}

% =====================================================================
%  30.3  The Matrix-Tree theorem.
%  Exports: sec:c32-matrix  def:c32-lap  thm:c32-matrix  lem:c32-cor
%  lst:c32-lap  lst:c32-moddet  fig:c32-lap  fig:c32-c5
%  WORKED examples only -- no solved problem here.  The Matrix-Tree
%  theorem gives spanning-tree count as any cofactor of the Laplacian;
%  a modular determinant is the whole algorithm for large counts.
%  Sanity checks against Cayley (K_n -> n^{n-2}), cycle C_n -> n, and
%  complete bipartite K_{a,b} -> a^{b-1} b^{a-1}.
% =====================================================================
\section{The Matrix-Tree Theorem}\label{sec:c32-matrix}

\begin{previouslybox}[Pr\"ufer handled the complete graph; now the graph is given]
\(\secref{c32-prufer}\) and \(\secref{c32-cayley}\) enumerate \emph{all}
labelled trees on \([n]\) or those satisfying a degree filter.  The input of
this section is different: an \emph{explicit} graph \(G=(V,E)\), and the
question is how many spanning trees \(G\) has.  When \(G=K_n\) the answer is
Cayley's \(n^{n-2}\); when \(G\) is sparser you cannot reduce to a Pr\"ufer
code filter -- you need a different tool.  The Matrix-Tree theorem
(\(\thmref{c32-matrix}\)) is that tool: a single determinant of a matrix you
can build in \(\bigO{n+m}\) gives the answer.
\end{previouslybox}

\begin{topicintro}
Attach to \(G\) its \term{Laplacian} \(L = D - A\) (\(\defref{c32-lap}\)), the
\(n\times n\) matrix with \(L_{uu}=\deg(u)\) on the diagonal and \(L_{uv} =
-\)(number of edges \(u\!-\!v\)) off.  Delete row and column \(i\) for any
fixed \(i\); the resulting \((n-1)\times(n-1)\) cofactor \(L^{(i)}\) has a
determinant equal to the number of spanning trees of \(G\)
(\(\thmref{c32-matrix}\)), independent of which \(i\) you chose.  For
unweighted graphs the count is an integer and can be huge; work modulo a
prime \(p\) and evaluate the determinant by Gaussian elimination in
\(\bigO{n^3}\) -- that is \(\coderef{c32-moddet}\).

The section is \emph{worked examples only}: no CF problem is solved here,
because the pure Matrix-Tree template almost always appears bundled with a
graph-construction twist (a problem you meet in Ch.~29 or Ch.~31).  What is
here is the mathematical engine itself and three sanity checks: cycle \(C_n\)
gives \(n\); complete bipartite \(K_{a,b}\) gives \(a^{b-1}b^{a-1}\);
complete graph \(K_n\) gives \(n^{n-2}\) recovering Cayley.

\textbf{Prerequisites.} Modular inverse (Ch.~6 background) via Fermat's little
theorem for a prime modulus; standard Gaussian elimination on a matrix.
Nothing else is assumed.

\textbf{After this section you can}
\begin{itemize}[nosep]
  \item build \(L(G)\) in \(\bigO{n+m}\) directly from an adjacency list; \hfill\lv{3}
  \item state the theorem and its dependence on the all-ones kernel of \(L\); \hfill\lv{4}
  \item evaluate the modular determinant of \(L^{(i)}\) via \(\coderef{c32-moddet}\); \hfill\lv{4}
  \item sanity-check the answer against Cayley and known families. \hfill\lv{4}
\end{itemize}
\end{topicintro}

% ---------------------------------------------------------------
\subsection{The Laplacian and the cofactor}
% ---------------------------------------------------------------

\begin{definition}{Laplacian matrix}{c32-lap}
Let \(G=(V,E)\) be a finite multigraph on \(V=[n]\) (parallel edges allowed;
self-loops ignored in \(L\)).  Its \term{Laplacian} is the \(n\times n\)
integer matrix
\[
  L(G)_{uv} \;=\; \begin{cases}
    \deg_G(u), & u=v,\\
    -(\text{number of edges }u\!-\!v), & u\ne v.
  \end{cases}
\]
Equivalently, \(L = D - A\) with \(D\) the degree diagonal and \(A\) the
adjacency matrix.
\end{definition}

\begin{lemma}{Kernel of the Laplacian}{c32-cor}
For every \(G\), \(L(G)\cdot \mathbf{1} = 0\) (each row of \(L\) sums to
zero).  If \(G\) is connected, the null-space of \(L(G)\) is exactly the
span of \(\mathbf{1}\), so \(\mathrm{rank}\,L(G) = n-1\); deleting any one
row and the corresponding column gives a nonsingular matrix.
\end{lemma}

\begin{proofbox}
Row \(u\) of \(L\) has diagonal \(\deg(u)\) and off-diagonals \(-1\) per edge
incident to \(u\); summing across the row contributes \(\deg(u) - \deg(u) =
0\).  If \(x\in\ker L\) then \(x^\top L x = \sum_{(u,v)\in E}(x_u-x_v)^2\);
this vanishes iff \(x_u=x_v\) on every edge, i.e.\ iff \(x\) is constant on
each connected component.  Connectivity forces \(x\propto\mathbf{1}\).
\(\Box\)
\end{proofbox}

\begin{theorem}{Matrix-Tree theorem (Kirchhoff)}{c32-matrix}
For a connected graph \(G\) on \([n]\), the number of its spanning trees
\(\tau(G)\) equals \(\det L^{(i)}(G)\) for any \(i\in[n]\), where \(L^{(i)}\)
is the \((n-1)\times(n-1)\) matrix obtained by deleting row and column \(i\)
from \(L(G)\).  For a disconnected \(G\), every \((n-1)\)-cofactor has
determinant zero (there are no spanning trees).
\end{theorem}

The theorem has four standard proofs -- via the Cauchy-Binet identity
applied to a reduced incidence matrix, via a minors-of-a-minor argument on
\(L\), via an electrical-network effective-resistance identity, and via a
combinatorial species bijection.  We give the Cauchy-Binet one, since it
builds directly on \(\lemref{c32-cor}\).

\begin{proofbox}
Pick an orientation of \(G\)'s edges, and let \(B\) be its \((n-1)\times m\)
reduced incidence matrix (row \(u\) is the incidence vector of edge \(u\)
minus the reference row \(i\)).  Then \(L^{(i)} = B\,B^\top\) (a standard
identity you can verify entrywise: the diagonal counts incident edges, the
off-diagonal subtracts the shared-edge indicator).  By Cauchy-Binet, \(\det
(BB^\top) = \sum_{S\subseteq[m], |S|=n-1} (\det B_S)^2\), where \(B_S\) is the
square submatrix on edge set \(S\).  A classical fact about \(B\) is that
\(\det B_S\in\{+1,-1\}\) iff \(S\) is the edge set of a spanning tree of
\(G\), and \(\det B_S = 0\) otherwise.  So each spanning tree contributes
exactly \(1\) to the sum, and non-tree subsets contribute \(0\).  Hence
\(\det L^{(i)} = \tau(G)\). \(\Box\)
\end{proofbox}

\begin{lstlisting}[caption={Build the Laplacian of a graph on n vertices with m edges, then delete row and column 0 (equivalent to any other by the theorem). Modular Gaussian elimination evaluates the determinant. ASCII only, identifiers without underscores}, label={lst:c32-lap}]
#include <bits/stdc++.h>
using namespace std;
typedef long long ll;
const ll MOD=1000000007LL;
ll modPow(ll b, ll e, ll m){
  ll r=1; b%=m;
  while(e>0){ if(e&1) r=(r*b)%m; b=(b*b)%m; e>>=1; }
  return r;
}
ll modInv(ll a){ return modPow(a, MOD-2, MOD); }
int main(){
  int n,m; if(scanf("%d %d",&n,&m)!=1) return 0;
  vector<vector<ll> > L(n, vector<ll>(n,0));
  for(int i=0;i<m;i++){
    int u,v; scanf("%d %d",&u,&v); u--; v--;
    L[u][u]++; L[v][v]++;
    L[u][v] = (L[u][v]-1+MOD)%MOD;
    L[v][u] = (L[v][u]-1+MOD)%MOD;
  }
  int sz=n-1;
  vector<vector<ll> > M(sz, vector<ll>(sz,0));
  for(int i=0;i<sz;i++) for(int j=0;j<sz;j++) M[i][j]=L[i+1][j+1];
  ll det=1;
  for(int col=0; col<sz; col++){
    int piv=-1;
    for(int r=col; r<sz; r++) if(M[r][col]%MOD!=0){ piv=r; break; }
    if(piv<0){ det=0; break; }
    if(piv!=col){ swap(M[piv], M[col]); det=(MOD-det)%MOD; }
    det=(det*M[col][col])%MOD;
    ll inv=modInv(M[col][col]);
    for(int r=col+1; r<sz; r++){
      ll factor = (M[r][col]*inv)%MOD;
      for(int c=col; c<sz; c++)
        M[r][c] = (M[r][c] - factor*M[col][c]%MOD + MOD)%MOD;
    }
  }
  printf("%lld\n", det);
  return 0;
}
\end{lstlisting}

\begin{lstlisting}[caption={Standalone modular determinant of an \(s\times s\) matrix under a prime \(p\), used repeatedly when the same Laplacian is queried for different minors. Gaussian elimination with row swap and multiplicative pivots; \(\bigO{s^3}\)}, label={lst:c32-moddet}]
ll modDet(vector<vector<ll> > M, ll p){
  int s=(int)M.size();
  ll det=1;
  for(int col=0; col<s; col++){
    int piv=-1;
    for(int r=col; r<s; r++) if(M[r][col]%p!=0){ piv=r; break; }
    if(piv<0) return 0;
    if(piv!=col){ swap(M[piv], M[col]); det=(p-det)%p; }
    det=(det*M[col][col])%p;
    ll inv=modPow(M[col][col], p-2, p);
    for(int r=col+1; r<s; r++){
      ll factor=(M[r][col]*inv)%p;
      for(int c=col; c<s; c++)
        M[r][c]=(M[r][c]-factor*M[col][c]%p+p)%p;
    }
  }
  return det;
}
\end{lstlisting}

% ---------------------------------------------------------------
\subsection{Worked examples: three sanity checks}
% ---------------------------------------------------------------

\subsection*{Worked 1: The cycle \(C_5\)}

\(C_5\) on labels \(1,2,3,4,5\) with edges \((1,2),(2,3),(3,4),(4,5),(5,1)\).
Its Laplacian is tridiagonal on the ring, degree \(2\) on the diagonal and
\(-1\) on the cyclic off-diagonals.  Delete row and column \(5\); the
remaining \(4\times 4\) matrix is the path Laplacian
\[
  L^{(5)}(C_5) = \begin{pmatrix} 2 &-1 & 0 & 0\\ -1 & 2 &-1 & 0\\ 0&-1&2&-1\\ 0&0&-1&2\end{pmatrix}.
\]
Its determinant is the Fibonacci-like sequence \(d_1=2,\, d_2=3,\, d_3=4,\,
d_4=5\) (each cofactor of a path Laplacian on \(k\) nodes with fixed
boundaries is \(k+1\)); the answer is \(5\).  And indeed \(C_5\) has
\emph{exactly} \(5\) spanning trees -- remove any one of the five edges.
The formula matches the combinatorics.

\subsection*{Worked 2: The complete bipartite graph \(K_{2,3}\)}

Partition \(A=\{1,2\}\), \(B=\{3,4,5\}\), every \(A\)-vertex joined to every
\(B\)-vertex; no intra-part edges.  Degrees: \(\deg(1)=\deg(2)=3,\)
\(\deg(3)=\deg(4)=\deg(5)=2\).  Delete row and column \(5\); the resulting
\(4\times 4\) cofactor on \(\{1,2,3,4\}\) is
\[
  L^{(5)} = \begin{pmatrix}3&0&-1&-1\\0&3&-1&-1\\-1&-1&2&0\\-1&-1&0&2\end{pmatrix}.
\]
Reduce: rows 3 and 4 differ by \((0,0,2,-2)\), giving a pivot at \((3,3)=2\);
rows 1 and 2 by \((3,0,-1,-1)\) versus \((0,3,-1,-1)\), independent.  A
short elimination produces \(\det = 12\).  The closed formula
\(a^{b-1}b^{a-1}=2^{2}\cdot 3^{1}=12\) agrees.

\subsection*{Worked 3: The complete graph \(K_4\)}

\(L(K_4)=4I - (J-I) = 4I - J + I = 5I - J\) with \(J\) the all-ones matrix.
Delete row/col 4 to get \(L^{(4)}=5I_3 - J_3 = \mathrm{diag}(5)-\mathrm{off\text{-}diag}(1)\).
Its eigenvalues are \(5-3=2\) (once, on \(\mathbf{1}\)) and \(5\) (twice, on
the orthogonal complement), so the determinant is \(2\cdot 5\cdot 5=50\)??  No
-- re-check.  \(L(K_4)\) has diagonal \(3\) (degree \(3\) in \(K_4\)), not
\(5\).  Correct matrix \(L^{(4)}\) on the top-left \(3\times 3\) block has
diagonal \(3\) and off-diagonal \(-1\), so \(L^{(4)}=4I_3 - J_3\).
Eigenvalues on the same split: \(4-3=1\) on \(\mathbf{1}\) and \(4\) twice on
the orthogonal complement, so \(\det = 1\cdot 4 \cdot 4 = 16\).  And indeed
Cayley \(n^{n-2}=4^2=16\).  The formula matches \(\thmref{c32-cayley}\).

\begin{figure}[H]
\centering
\begin{tikzpicture}[diagram, scale=0.9, transform shape,
   tnode/.style={circle, draw=primaryblue!70, fill=primaryblue!8, line width=0.7pt,
                 minimum size=6mm, font=\scriptsize\bfseries, inner sep=1pt}]
  \node[tnode] (a) at (0,1) {1};
  \node[tnode] (b) at (1.2,1) {2};
  \node[tnode] (c) at (0,0) {3};
  \node[tnode] (d) at (1.2,0) {4};
  \node[tnode] (e) at (0.6,-1) {5};
  \draw (a)--(c)--(e)--(d)--(b)--(a);
  \draw (a)--(b); \draw (c)--(d);
  \node[annot, font=\scriptsize, anchor=north] at (0.6,-2.4)
    {\(C_5\) on five labels: delete any row/col pair and the \(4\times 4\) cofactor has determinant \(5\)};
\end{tikzpicture}
\caption{The five-cycle \(C_5\); \(\tau=5\) because a spanning tree is
exactly the graph minus one of the five edges.  Worked check of
\(\thmref{c32-matrix}\) against the combinatorial count.}
\label{fig:c32-c5}
\end{figure}

\begin{insightbox}[What Matrix-Tree buys you that Pr\"ufer does not]
Pr\"ufer's bijection is a closed form on \(\mathcal{T}_n\); it is unbeatable
when the answer is a sequence enumeration.  Matrix-Tree is the \emph{generic
engine} -- hand it any explicit graph and it returns the count.  The two
tools are not rivals; the problems that mix them (given a graph, count its
spanning trees \emph{with a degree constraint}) are exactly the CF 2400+ tier
where you must decide which tool handles which part of the constraint.
Section \(\secref{c32-boundary}\) makes the split precise, and the Summit
\(\probref{c32-summit}\) uses the completion formula from
\(\secref{c32-cayley}\) rather than the determinant because the graph there
is a \emph{given tree}, not an arbitrary one.
\end{insightbox}

\begin{notebox}[Modular pitfalls]
\textbf{Modular pitfalls.}  The determinant algorithm assumes the modulus is
prime (or at least that every pivot is invertible).  Use \(p=10^9+7\) or
\(p=998244353\); avoid composite moduli unless you switch to Bareiss' exact
division (slower but works over \(\mathbb{Z}\)).  Row swaps flip the sign,
handled by \((\mathrm{MOD}-\det)\bmod\mathrm{MOD}\).  Zero pivots on a
nonzero determinant are handled by picking a row below to swap in.
\end{notebox}

% =====================================================================
%  30.4  Linearity of expectation on a rooted tree.
%  Exports: sec:c32-expect  def:c32-ind  lem:c32-indexp  thm:c32-sum
%  prob:c32-game  fig:c32-game  lst:c32-game  lst:c32-depth
%  Solved CF 280C "Game on Tree" (2200, L4), live-verified against the CF
%  problem page (title "C. Game on Tree", 1s TL, 256 MB, n<=1e5, answer
%  tolerance 1e-6).  Two official samples: n=2 edge 1-2 -> 1.5; n=3 star
%  1-2, 1-3 -> 2.0.  Formula E = sum_v 1/depth(v), depth counted in
%  vertices on the root-to-v path so depth(root)=1.
% =====================================================================
\section{Linearity of Expectation on a Rooted Tree}\label{sec:c32-expect}

\begin{previouslybox}[Counting becomes expectation with one small twist]
\(\secref{c32-prufer}\) through \(\secref{c32-matrix}\) produce an integer
answer (often modulo a prime).  This section produces a real: an expectation.
The trick is the same trick a probabilist uses everywhere -- split a random
variable into indicators, and swap the sum with the expectation.  On trees
the split is naturally \emph{per vertex}, and the resulting probability is a
reciprocal of a depth, so \(\secref{c32-expect}\) hands you an \(O(n)\) DFS
where a naive recursion on the whole process would be exponential.
\end{previouslybox}

\begin{topicintro}
Let \(T\) be a rooted tree with root \(r=1\).  A random process on \(T\) has
a state space that grows exponentially with \(n\); writing down a recursion
for \(\mathbb{E}[\text{steps}]\) is hopeless.  But expectation is \emph{linear
without independence} (\(\lemref{c32-indexp}\)): for any random variables
\(X_1,\dots,X_n\) (dependent or not), \(\mathbb{E}[\sum X_i]=\sum\mathbb{E}[X_i]\).
Decompose "number of steps" as \(\sum_v X_v\), where \(X_v=1\) iff \(v\) is
\emph{directly picked} by the process, and each \(\mathbb{E}[X_v] =
\Pr[X_v=1]\) becomes a small rational (\(\thmref{c32-sum}\)).  On the CF 280C
process -- repeatedly choose a remaining vertex uniformly and delete its
subtree -- the answer collapses to \(\sum_v 1/\mathrm{depth}(v)\), a single
DFS.

CF 280C (\(\probref{c32-game}\)) is exactly that computation.  It is a
\(\lv{4}\) problem not because the mathematics is hard but because the
\emph{decomposition} is not obvious -- most first attempts try to write a
DP over the process, which is exponential.

\textbf{Prerequisites.} Ch.~12 (\chref{c09}) for the depth array; basic
probability (uniform distribution over a finite set).  Nothing else.

\textbf{After this section you can}
\begin{itemize}[nosep]
  \item state linearity of expectation without an independence hypothesis; \hfill\lv{3}
  \item decompose a step-count into per-vertex indicators; \hfill\lv{4}
  \item compute \(\Pr[v\text{ is picked}]\) via a symmetry argument; \hfill\lv{4}
  \item solve CF 280C in one DFS. \hfill\lv{4}
\end{itemize}
\end{topicintro}

% ---------------------------------------------------------------
\subsection{Linearity without independence}
% ---------------------------------------------------------------

\begin{definition}{Indicator variable of an event}{c32-ind}
For an event \(A\) in a probability space \(\Omega\), its \term{indicator}
\(\mathbf{1}_A:\Omega\to\{0,1\}\) is \(\mathbf{1}_A(\omega)=1\) iff
\(A\) occurs at \(\omega\).  \(\mathbb{E}[\mathbf{1}_A]=\Pr[A]\).
\end{definition}

\begin{lemma}{Linearity of expectation}{c32-indexp}
For random variables \(X_1,\dots,X_n\) on a common probability space, not
assumed independent, \(\mathbb{E}\!\left[\sum_i X_i\right]=\sum_i
\mathbb{E}[X_i]\).
\end{lemma}

\begin{proofbox}
Follows from \(\mathbb{E}[X]=\sum_{\omega} P(\omega) X(\omega)\) and the
distributive law: swap the order of the two finite sums.  No independence
assumption is used. \(\Box\)
\end{proofbox}

\begin{theorem}{Depth-reciprocal sum for CF 280C}{c32-sum}
Let \(T\) be a rooted tree on \([n]\) with root \(1\).  Consider the process
that repeatedly picks a uniformly random remaining vertex and deletes its
subtree, until the tree is empty.  The expected number of steps is
\[
  \mathbb{E}[\text{steps}] \;=\; \sum_{v=1}^{n} \frac{1}{\mathrm{depth}(v)},
\]
where \(\mathrm{depth}(v)\) is the number of vertices on the root-\(v\)
path, so \(\mathrm{depth}(1)=1\).
\end{theorem}

\begin{proofbox}
The process is equivalent to assigning each vertex \(v\) an independent
uniform random priority \(\pi_v\in[0,1]\), then processing vertices in
increasing priority order and marking a step for every vertex whose
ancestors on the root-path have all been processed (i.e., were the "next
chosen" at their own moment).  A vertex \(v\) contributes one step iff
\(\pi_v\) is the minimum on the root-\(v\) path; by symmetry over the
\(\mathrm{depth}(v)\) path vertices, \(\Pr[\pi_v\text{ is the min}]=1/\mathrm{depth}(v)\).

Now write \(X_v = \mathbf{1}\{v\text{ contributes a step}\}\).  Total steps
\(N = \sum_v X_v\).  By \(\lemref{c32-indexp}\),
\(\mathbb{E}[N] = \sum_v \Pr[X_v=1] = \sum_v 1/\mathrm{depth}(v)\). \(\Box\)
\end{proofbox}

\begin{lstlisting}[caption={Root a tree at 1 and compute depth via a single DFS or BFS; the expectation is the sum of reciprocal depths. Iterative DFS with an explicit stack avoids recursion-depth issues on n up to 1e5. ASCII only, identifiers without underscores}, label={lst:c32-depth}]
#include <bits/stdc++.h>
using namespace std;
int main(){
  int n; if(scanf("%d",&n)!=1) return 0;
  vector<vector<int> > adj(n+1);
  for(int i=0;i<n-1;i++){
    int u,v; scanf("%d %d",&u,&v);
    adj[u].push_back(v); adj[v].push_back(u);
  }
  vector<int> depth(n+1,0);
  depth[1]=1;
  vector<int> st; st.push_back(1);
  vector<int> par(n+1,0);
  while(!st.empty()){
    int u=st.back(); st.pop_back();
    for(int v : adj[u]) if(v!=par[u]){
      par[v]=u; depth[v]=depth[u]+1; st.push_back(v);
    }
  }
  double E=0.0;
  for(int v=1; v<=n; v++) E += 1.0/depth[v];
  printf("%.20f\n", E);
  return 0;
}
\end{lstlisting}

% ---------------------------------------------------------------
\subsection{Solved: CF 280C -- Game on Tree (2200, L4)}
% ---------------------------------------------------------------

\solvedproblem{c32-game}{Game on Tree}{CF 280C}{https://codeforces.com/problemset/problem/280/C}{CF 2200 L4}{Linearity of expectation (\(\lemref{c32-indexp}\)); per-vertex indicator decomposition; one DFS for depth (\(\thmref{c32-sum}\), \(\coderef{c32-depth}\))}

\begin{problemstatement}[CF 280C, verified against the live CF problem page]
A rooted tree on \(n\) vertices has root \(1\); the next \(n-1\) lines
describe edges as pairs \((a_i,b_i)\) with \(1\le a_i,b_i\le n\) and
\(a_i\ne b_i\), and the graph is guaranteed to be a tree.  A one-player game
runs as follows: while at least one vertex remains, pick a uniformly random
vertex among those remaining and delete its whole rooted subtree (the vertex
and every descendant of it in the residual forest).  The game ends when no
vertices remain.  Report the expected number of steps as a real number to
absolute or relative error \(\le 10^{-6}\).  Constraint: \(1\le n\le
10^5\).  Samples: input \(2,1\;2\) outputs \(1.5\); input \(3,1\;2,1\;3\)
outputs \(2.0\).
\end{problemstatement}

\begin{figure}[H]
\centering
\begin{tikzpicture}[diagram, scale=0.9, transform shape,
   tnode/.style={circle, draw=primaryblue!70, fill=primaryblue!8, line width=0.7pt,
                 minimum size=6mm, font=\scriptsize\bfseries, inner sep=1pt}]
  \node[tnode] (r) at (0,1.2) {1};
  \node[tnode] (a) at (-1.2,-0.2) {2};
  \node[tnode] (b) at (1.2,-0.2) {3};
  \draw[primaryblue!70, line width=1pt] (r)--(a) (r)--(b);
  \node[annot, font=\scriptsize, anchor=north] at (0,-1.4)
    {sample 2: depths \((1,2,2)\), \(E = 1 + 1/2 + 1/2 = 2\)};
\end{tikzpicture}
\caption{Sample 2 of CF 280C.  Root \(1\) has depth \(1\); vertices \(2\) and
\(3\) have depth \(2\).  The depth-reciprocal sum gives \(E=2\), matching the
published sample output \(2.00000000000000000000\).}
\label{fig:c32-game}
\end{figure}

\begin{attemptgate}[Attempt it yourself first -- \lv{4}, give yourself 25 minutes]
Before coding: can you write a recursion on "expected steps from a residual
forest \(F\)"?  What goes wrong, and what decomposition avoids it?
\peek{A recursion on \(F\) has \(2^n\) states; memoising is not an option.
Instead, ask: what is one \emph{step} made of?  It picks exactly one vertex,
so the total number of steps equals the number of vertices DIRECTLY picked
(never covered by an ancestor).}
\nudge{Call that event \(A_v\).  The total steps \(N = \sum_v
\mathbf{1}_{A_v}\).  Linearity of expectation (\(\lemref{c32-indexp}\)) gives
\(\mathbb{E}[N] = \sum_v \Pr[A_v]\), and the sum runs over \(n\) terms with
no independence required.}
\fullhint{The process on \(T\) is equivalent to assigning each vertex an
i.i.d.\ continuous uniform priority and processing vertices in increasing
priority order; a vertex \(v\) is "directly picked" iff its priority is the
minimum on the root-\(v\) path.  By symmetry over the path's
\(\mathrm{depth}(v)\) labels, this has probability \(1/\mathrm{depth}(v)\).
Compute depths with one DFS (\(\coderef{c32-depth}\)) and sum the
reciprocals; \(\bigO{n}\) time, \(\bigO{n}\) space.}
\end{attemptgate}

\solstep{Understanding}
CF 280C looks like a dynamic-programming-on-subtrees problem because the
input is a rooted tree, but the state is a \emph{set of remaining vertices},
which explodes exponentially.  The rescue is \(\lemref{c32-indexp}\): you
never need to reason about the whole process, only about whether each single
vertex is picked directly.  The indicator decomposition collapses the
exponential state space to \(n\) independent questions, each of which has
the same answer for the same-depth vertices.  What remains is a formula,
and the formula is one DFS.

The priority-ordering equivalence spelled out in the full hint above is standard: drawing
each vertex uniformly at random from \([0,1]\) and taking them in increasing
order is a uniformly random permutation of the vertices.  The process
"pick a uniformly random remaining vertex at each step" is a uniformly
random permutation as well, and it deletes a vertex's subtree at the
vertex's own moment in the permutation.

\solstep{Key Observations}
\begin{observation}
Every step picks exactly one vertex that survives until that step, and
deletes a set of vertices (its subtree).  The vertex \emph{doing} the
deletion at step \(t\) is called the "picked" vertex; every other deleted
vertex is "covered" by an ancestor's pick.  Total steps \(=\) total picked
vertices \(=\sum_v \mathbf{1}_{A_v}\).
\end{observation}
\begin{observation}
\(A_v\) holds iff \(v\) appears before any of its ancestors in the random
permutation.  Equivalently, among the \(\mathrm{depth}(v)\) vertices on the
root-to-\(v\) path, \(v\) is the earliest.  By symmetry of a uniformly
random permutation, each of these \(\mathrm{depth}(v)\) labels is equally
likely to be the earliest, so \(\Pr[A_v]=1/\mathrm{depth}(v)\).
\end{observation}
\begin{observation}
No independence between \(A_u\) and \(A_v\) is required.  Linearity of
expectation (\(\lemref{c32-indexp}\)) sums \(n\) rationals regardless of the
dependence structure among the indicators.
\end{observation}

\solstep{Think It With Me}
\begin{thinkbox}
\deadend{Recursion on subtrees: \(\mathbb{E}[T] = 1 + \frac{1}{|V(T)|}
\sum_v \mathbb{E}[T\setminus \mathrm{subtree}(v)]\)}{The state space is the
set of residual forests reachable by subtree deletions; there are
exponentially many, and the recursion has no memoisable subproblem.}
\deadend{Try a symmetry argument over "expected depth"}{That measures only a single vertex's path contribution, not the sum of step indicators, so it undercounts the total expectation.}
\deadend{Write a DP with state "root picked or not"}{That flag captures whether the root has fired, but the residual forest below the root still has exponentially many configurations once you recurse.}
\aha{Split into indicators, one per vertex.  Linearity applies with no
independence, and each indicator's probability is \(1/\mathrm{depth}(v)\)
by symmetry of a uniform permutation.}
\end{thinkbox}

\solstep{The Approach Ladder}
\begin{center}\renewcommand{\arraystretch}{1.25}\setlength{\tabcolsep}{3pt}
\begin{tabular}{@{}c p{0.26\linewidth} p{0.30\linewidth} l c@{}}
\toprule
\textbf{\#} & \textbf{Approach} & \textbf{Idea} & \textbf{Time} & \textbf{Verdict}\\
\midrule
1 & \lv{2}\ full recursion & exponential state space & \(\bigO{2^n}\) & \tle\\
2 & \lv{3}\ DP with "root picked" flag & still exponential residual forest & \(\bigO{2^n}\) & \tle\\
3 & \lv{4}\ \textbf{indicator + linearity} & per-vertex symmetry & \(\bigO{n}\) & \ac\\
\bottomrule
\end{tabular}\end{center}
Rungs 1 and 2 are natural but doomed.  Rung 3 is one DFS and one sum.

\solstep{Solution}
Root the tree at vertex \(1\) (which the statement already fixes).  Run one
DFS to fill \(\mathrm{depth}[v]\) for every \(v\) (\(\coderef{c32-depth}\));
the recursion is \(O(n)\) total.  Then compute \(\sum_v 1/\mathrm{depth}[v]\)
as a \code{double} and print with sufficient precision (\texttt{\%.20f} is a
safe default for a \(10^{-6}\) error budget).

\solstep{Implementation}
\begin{lstlisting}[caption={CF 280C. Iterative DFS for depth, sum reciprocal depths, print with wide precision. The iterative stack avoids recursion-depth issues at n=1e5. ASCII only, identifiers without underscores}, label={lst:c32-game}]
#include <bits/stdc++.h>
using namespace std;
int main(){
  int n; if(scanf("%d",&n)!=1) return 0;
  vector<vector<int> > adj(n+1);
  for(int i=0;i<n-1;i++){
    int u,v; scanf("%d %d",&u,&v);
    adj[u].push_back(v); adj[v].push_back(u);
  }
  vector<int> depth(n+1,0), par(n+1,0);
  vector<int> st; st.push_back(1); depth[1]=1;
  while(!st.empty()){
    int u=st.back(); st.pop_back();
    for(int v : adj[u]) if(v!=par[u]){
      par[v]=u; depth[v]=depth[u]+1; st.push_back(v);
    }
  }
  double E=0.0;
  for(int v=1; v<=n; v++) E += 1.0/depth[v];
  printf("%.20f\n", E);
  return 0;
}
\end{lstlisting}

\solstep{Proof of Correctness}
\begin{proofbox}
\textbf{(1)} The process is equivalent to a uniformly random permutation of
\([n]\): at each step pick the surviving vertex with the smallest position
in the permutation, and delete its subtree in the surviving forest.  This
generates the same distribution over "which vertex is picked at each step"
as the original process, because at every step the choice among the
survivors is uniform by the permutation's symmetry.

\textbf{(2)} A vertex \(v\) is picked by this permutation-protocol iff \(v\)
has the smallest permutation position among all vertices on the root-\(v\)
path in the original tree.  Reason: if an ancestor \(a\) of \(v\) has a
smaller position, then \(a\) is picked first and its subtree-deletion
removes \(v\), so \(v\) is never directly picked.  Conversely if \(v\) has
the smallest position on the path, then no ancestor of \(v\) has fired
before \(v\) does, so \(v\) survives until its own turn.

\textbf{(3)} By symmetry of a uniformly random permutation restricted to
any fixed set of \(\mathrm{depth}(v)\) positions, the smallest of them is
uniformly distributed over those \(\mathrm{depth}(v)\) vertices, so
\(\Pr[A_v] = 1/\mathrm{depth}(v)\).

\textbf{(4)} \(N=\sum_v \mathbf{1}_{A_v}\), so by \(\lemref{c32-indexp}\)
(linearity without independence) \(\mathbb{E}[N]=\sum_v 1/\mathrm{depth}(v)\),
which is what \(\coderef{c32-game}\) computes. \(\Box\)
\end{proofbox}

\solstep{Trace on the Official Samples}
\begin{tracebox}[Sample 1: \, \(n=2\), edge \(1\!-\!2\)]
Depths \(\mathrm{depth}[1]=1\), \(\mathrm{depth}[2]=2\).  Sum \(=1+1/2=
1.5\).  The published sample output is \(1.50000000000000000000\).
Direct check: two runs of the game -- either pick \(1\) first (probability
\(1/2\), \(1\) step; the tree empties immediately) or pick \(2\) first
(probability \(1/2\), \(1\) step deletes \(2\), then one more step deletes
residual \(1\), so \(2\) steps).  \(\mathbb{E}=1/2\cdot 1 + 1/2\cdot 2 =
1.5\).  Matches.
\end{tracebox}
\begin{tracebox}[Sample 2: \, \(n=3\), edges \(1\!-\!2\) and \(1\!-\!3\)]
Depths \((1,2,2)\).  Sum \(=1+1/2+1/2=2\).  Published sample output
\(2.00000000000000000000\).  Direct: pick root \(1\) (probability \(1/3\))
ends in \(1\) step; else pick \(2\) or \(3\) first (probability \(2/3\)),
which deletes one leaf; the residual tree is an edge on \(\{1,3\}\) or
\(\{1,2\}\), expected \(1.5\) more steps; total \(1/3\cdot 1 + 2/3\cdot
(1 + 1.5) = 1/3 + 5/3 = 2\).  Matches.
\end{tracebox}

\solstep{Complexity}
\begin{complexitybox}
Building the adjacency list: \(\bigO{n}\).  One iterative DFS fills
\(\mathrm{depth}[\cdot]\): \(\bigO{n}\).  Sum of reciprocals: \(\bigO{n}\).
Total \(\bigO{n}\) time, \(\bigO{n}\) space.  For \(n=10^5\) this is
comfortably under the 1-second limit.  The \code{double} accumulator has a
round-off error of \(\bigO{n\epsilon_{\mathrm{machine}}}\approx 10^{-11}\)
for \(n=10^5\), far inside the \(10^{-6}\) tolerance.
\end{complexitybox}

\begin{hardeningbox}
\begin{itemize}[nosep,leftmargin=1.4em]
\item \textbf{Iterative DFS.}  \(n\le 10^5\), and a chain of that length
      recurses \(10^5\) frames deep -- an easy stack overflow with default
      settings.  The explicit-stack form in \(\coderef{c32-game}\) is
      safer and equally short.
\item \textbf{Depth starts at 1.}  A vertex at BFS level 0 (the root) is at
      path-length \(\mathrm{depth}=1\) (root itself).  Off-by-one here
      changes the answer by exactly \(1\) on the sample \(n=2\) test.
\item \textbf{Print format.}  \texttt{\%.20f} is overkill but safe; do not
      switch to \texttt{\%g}, which may print scientific notation and can
      drop significant digits under some libraries.
\item \textbf{Adjacency vector of vectors.}  Fine for \(n\le 10^5\); for
      tighter memory budgets use a flat edge list with head/next arrays.
\end{itemize}
\end{hardeningbox}

\begin{recallbox}[30.4 in one breath]
\emph{Every step picks exactly one vertex; total steps is \(\sum_v
\mathbf{1}_{A_v}\).  Linearity of expectation (\(\lemref{c32-indexp}\))
sums the \(\Pr[A_v]\) with no independence assumption.  Symmetry of a
uniform permutation on a fixed path gives \(\Pr[A_v]=1/\mathrm{depth}(v)\).
One DFS, one sum, done (\(\thmref{c32-sum}\), \(\coderef{c32-game}\)).}
\selftest{
  \item State linearity of expectation and name one hypothesis it does
        \emph{not} require.
  \item Explain why the "random permutation" reformulation of CF 280C
        preserves the process's distribution.
  \item Compute the answer by hand for a path \(1\!-\!2\!-\!3\) rooted at \(1\).
        \emph{(Depths \((1,2,3)\), sum \(1+1/2+1/3 = 11/6\).)}
}
\end{recallbox}

% =====================================================================
%  30.5  The method boundary: labelled vs unlabelled, and where each
%  tool of this chapter stops applying.
%  Exports: sec:c32-boundary  def:c32-orbit  thm:c32-orbit  def:c32-pick
%  fig:c32-orbit
%  WORKED examples only.  Hands off to Ch.28 (shapes), Ch.26 (virtual
%  trees), Ch.31 (MST).  The reader leaves Ch.30 with the split "count
%  sequences / take a determinant / split into indicators / hand off".
% =====================================================================
\section{The Method Boundary: Labelled vs.\ Unlabelled}\label{sec:c32-boundary}

\begin{previouslybox}[Ch.28 compared shapes; Ch.30 counted labelings]
Chapter~28 (\chref{c22}) took two trees and decided whether they were the
\emph{same shape}, running AHU or a tree hash to canonicalise.  This chapter
counts \emph{labelled} trees, where two assignments of the same shape to the
vertex set \([n]\) are different objects.  The two views are reconciled by a
single group-action identity: the size of the orbit of a labelled tree under
\(S_n\) is \(n!/|\mathrm{aut}(T)|\).  Section \(\secref{c32-boundary}\) makes that
precise, hands the unlabelled case back to Ch.~28, and lists three nearby
topics -- virtual trees (Ch.~26 \chref{c30}) and MST (Ch.~31 \chref{c17})
among them -- that are \emph{not} this chapter's territory.
\end{previouslybox}

\begin{topicintro}
Two questions that sound identical and are not: "How many trees are there on
\(n\) vertices?" and "How many trees are there with \(n\) vertices?".  The
first is Cayley's \(n^{n-2}\) (labelled: which subset of \([n]\times[n]\) is
the edge set).  The second is the count of \emph{unlabelled tree shapes}
(A000055 in OEIS); the answer is \(1, 1, 1, 2, 3, 6, 11, 23, 47, 106,\dots\)
and has no closed form -- it is computed by P\'olya enumeration or by the
rooted-shape recursion of Ch.~28.  Deciding which question a statement asks
is the single most important reading skill of this section.

The bridge between the two is the orbit-stabiliser theorem
(\(\thmref{c32-orbit}\)): \(S_n\) acts on \(\mathcal{T}_n\) by relabelling,
and each unlabelled shape \(S\) contributes \(n!/|\mathrm{aut}(S)|\)
labelled trees to the total.  Cayley's number is therefore \(\sum_{S\text{
shape on }n} n!/|\mathrm{aut}(S)|\).  This identity is why the unlabelled
enumeration is harder -- \(|\mathrm{aut}(S)|\) depends on the shape's
symmetry structure, which Ch.~28 is precisely the machinery to compute.

The rest of this section is a \emph{pick list}: four questions you will meet
that are \emph{not} Chapter~30's to answer, with a one-line pointer to the
chapter that owns each.  The chapter ends with the Summit CF 917D
(\(\probref{c32-summit}\)), which is a genuine Ch.~30 question mark,
unlabelled on the surface but labelled underneath.

\textbf{Prerequisites.} Ch.~28 (\chref{c22}) for \(|\mathrm{aut}(T)|\) via
AHU; Ch.~26 (\chref{c30}) for virtual trees; Ch.~31 (\chref{c17}) for MST.

\textbf{After this section you can}
\begin{itemize}[nosep]
  \item distinguish labelled from unlabelled counting in a problem statement; \hfill\lv{4}
  \item apply orbit-stabiliser to convert between the two views; \hfill\lv{5}
  \item name the three nearest topics this chapter does \emph{not} own and cite their chapters. \hfill\lv{4}
\end{itemize}
\end{topicintro}

% ---------------------------------------------------------------
\subsection{The labelled vs.\ unlabelled split}
% ---------------------------------------------------------------

\begin{definition}{Automorphism of a tree}{c32-orbit}
An \term{automorphism} of a labelled tree \(T\in\mathcal{T}_n\) is a
permutation \(\sigma\in S_n\) with \(\{u,v\}\in E(T)\iff \{\sigma(u),
\sigma(v)\}\in E(T)\).  The automorphism group \(\mathrm{aut}(T)\) is the
set of such \(\sigma\); its size measures the shape's internal symmetry.
\end{definition}

\begin{theorem}{Labelled count via unlabelled shapes}{c32-orbit}
Let \(\mathcal{S}_n\) be the set of unlabelled tree shapes on \(n\) vertices.
Then \(|\mathcal{T}_n|=\sum_{S\in\mathcal{S}_n} n!/|\mathrm{aut}(S)|\).
\end{theorem}

\begin{proofbox}
The symmetric group \(S_n\) acts on \(\mathcal{T}_n\) by relabelling.  Two
labelled trees \(T,T'\) are in the same orbit iff they share an unlabelled
shape.  By orbit-stabiliser, the size of the orbit containing \(T\) is
\([S_n:\mathrm{Stab}(T)]=n!/|\mathrm{Stab}(T)|\), and \(\mathrm{Stab}(T)=
\mathrm{aut}(T)\).  Summing orbit sizes over the set of shapes gives the
total \(|\mathcal{T}_n|\). \(\Box\)
\end{proofbox}

\subsection*{Worked: reconciling \(n=4\)}

Ch.~28 has \(2\) unlabelled tree shapes on \(4\) vertices: the path \(P_4\)
and the star \(K_{1,3}\).  \(|\mathrm{aut}(P_4)|=2\) (the identity and the
end-to-end flip), \(|\mathrm{aut}(K_{1,3})|=6\) (all permutations of the
three leaves).  Orbit sizes: \(4!/2 = 12\) labelled copies of \(P_4\), and
\(4!/6 = 4\) labelled copies of \(K_{1,3}\).  Sum \(=16\), matching Cayley
\(4^{4-2}=16\).  Both views agree.

\begin{figure}[H]
\centering
\begin{tikzpicture}[diagram, scale=0.9, transform shape,
   tnode/.style={circle, draw=primaryblue!70, fill=primaryblue!8, line width=0.7pt,
                 minimum size=6mm, font=\scriptsize\bfseries, inner sep=1pt}]
  \node[tnode] (p1) at (0,0) {1};
  \node[tnode] (p2) at (1,0) {2};
  \node[tnode] (p3) at (2,0) {3};
  \node[tnode] (p4) at (3,0) {4};
  \draw (p1)--(p2)--(p3)--(p4);
  \node[annot, font=\scriptsize] at (1.5,-0.9) {\(P_4\): aut size 2, orbit 12};
  \node[tnode] (c) at (6,0) {1};
  \node[tnode] (l1) at (5,1) {2};
  \node[tnode] (l2) at (7,1) {3};
  \node[tnode] (l3) at (7,-1) {4};
  \draw (c)--(l1) (c)--(l2) (c)--(l3);
  \node[annot, font=\scriptsize] at (6,-2) {\(K_{1,3}\): aut size 6, orbit 4};
\end{tikzpicture}
\caption{The two unlabelled tree shapes on \(4\) vertices.  Their orbit
sizes add to Cayley's \(16\).}
\label{fig:c32-orbit}
\end{figure}

% ---------------------------------------------------------------
\subsection{The pick list: what this chapter does not own}
% ---------------------------------------------------------------

\begin{definition}{Method-boundary pick list}{c32-pick}
Read the statement and pick:
\begin{itemize}[nosep]
  \item \textbf{Labelled count, closed form possible} -- Pr\"ufer
        (\(\secref{c32-prufer}\)) or Cayley variants
        (\(\secref{c32-cayley}\)).
  \item \textbf{Labelled count inside a \emph{given} graph} -- Matrix-Tree
        (\(\secref{c32-matrix}\)).
  \item \textbf{Expected value of a random process on a tree} -- linearity
        of expectation (\(\secref{c32-expect}\)).
  \item \textbf{Unlabelled shape census, automorphism size, tree isomorphism}
        -- Ch.~28 (\chref{c22}); use \(\thmref{c32-orbit}\) to reconcile
        against a labelled count if the problem mixes the two views.
  \item \textbf{Minimum spanning tree of a weighted graph} -- Ch.~31
        (\chref{c17}); the "tree" is a \emph{specific} one, not a count.
  \item \textbf{Tree over a marked vertex subset, distance queries} --
        Ch.~26 (\chref{c30}) virtual trees; the object is a small tree
        built per query, not a family.
\end{itemize}
\end{definition}

\begin{insightbox}[The one question you should be able to answer]
Given a problem statement, is the answer \emph{a number mod a prime}, a
\emph{real number}, or \emph{a specific tree}?  A number mod a prime means
you are counting labelled or spanning trees (Ch.~30).  A real number means
expectation (Ch.~30.4 is your template).  A specific tree means construction,
which is Ch.~26 or Ch.~31 territory.  This three-way split is the whole of
\(\secref{c32-boundary}\); the rest of the book elaborates.
\end{insightbox}

\begin{notebox}[A cautionary "count unlabelled trees" example]
"How many unlabelled tree shapes on \(n=8\) vertices?" is a Ch.~28 question;
the answer is \(23\).  No Ch.~30 tool applies -- Pr\"ufer is a bijection on
\emph{labelled} trees and Matrix-Tree counts \emph{labelled} spanning trees
of a specific graph.  Trying to "divide Cayley's \(8^6\) by \(8!\)" is a
classic wrong move; the orbit sizes vary and you cannot factor out a
uniform \(n!\).  Ch.~28's AHU or P\'olya machinery handles the enumeration.
\end{notebox}

\begin{recallbox}[30.5 in one breath]
\emph{Labelled trees form an \(S_n\)-set; unlabelled shapes are its orbits,
and \(\thmref{c32-orbit}\) is the bridge.  Cayley counts labels, AHU counts
shapes, and mixing the two without the orbit correction is the classic
error.  Virtual trees (Ch.~26 \chref{c30}) and MST (Ch.~31 \chref{c17}) are
one step outside this chapter's territory -- leave the answer space and you
leave Ch.~30.}
\selftest{
  \item Reconcile \(|\mathcal{T}_4| = 16\) via the two shapes on \(4\)
        vertices.
  \item Explain why dividing Cayley by \(n!\) is wrong for unlabelled
        enumeration.
  \item Name one Ch.~26 problem that is \emph{not} a Ch.~30 problem.
}
\end{recallbox}

% =====================================================================
%  Chapter 30 station file -- 14 chaptersections + Summit + lastword.
%    1 Summary                8 Recommended Study Order
%    2 Master Decision Flow   9 Difficulty Progression
%    3 Common Mistakes       10 Reduction Gallery
%    4 WA Clinic             11 Mixed Problem Set (10 practice items)
%    5 Quick Reference       12 Summit -- CF 917D Stranger Trees
%    6 Write-From-Memory     13 Tool Kit Inventory
%    7 Complete Checklist    14 Related Topics
%  Summit live-verified: title "D. Stranger Trees", 1s TL, 256 MB, n<=100,
%  mod 10^9+7.  Samples: star on 3 -> "0 2 1"; path on 4 -> "1 7 7 1";
%  star on 4 -> "0 9 6 1".  All three hand-replayed in the Trace section.
% =====================================================================

\chaptersection{Chapter Summary}

Chapter~30 is the enumerative counterweight to Chapters 12--29.  Where those
chapters \emph{operate} on a tree this one \emph{counts} trees.  Four tools
carry the chapter:

\begin{itemize}[nosep]
  \item the \textbf{Pr\"ufer bijection} (\(\secref{c32-prufer}\)):
        \(\mathcal{T}_n\leftrightarrow[n]^{n-2}\), with the code's multiplicity
        of label \(v\) equal to \(\deg_T(v)-1\);
  \item \textbf{Cayley's formula and its variants} (\(\secref{c32-cayley}\)):
        \(n^{n-2}\), rooted \(n^{n-1}\), degree multinomial \((n-2)!/\prod(d_i-1)!\),
        forest-completion \(n^{k-2}\prod_i s_i\);
  \item the \textbf{Matrix-Tree theorem} (\(\secref{c32-matrix}\)): spanning
        trees of a given graph \(G\) equal any cofactor of \(L(G)=D-A\); a
        modular determinant is the algorithm;
  \item \textbf{linearity of expectation} (\(\secref{c32-expect}\)): per-vertex
        indicator decomposition, with the CF 280C process collapsing to
        \(\sum_v 1/\mathrm{depth}(v)\).
\end{itemize}

\(\secref{c32-boundary}\) closes the loop by drawing the line between
\emph{labelled} counts (this chapter) and \emph{unlabelled} shapes
(Ch.~28 \chref{c22}), via the orbit-stabiliser bridge
\(\thmref{c32-orbit}\).  The Summit \(\probref{c32-summit}\) is a genuine
labelled-count question mark whose answer uses the forest-completion formula
as its workhorse and a binomial inversion to close the "exactly \(k\) shared
edges" filter.

\chaptersection{Master Decision Flowchart}

\begin{center}
\begin{tikzpicture}[diagram, scale=0.85, transform shape,
   q/.style={rectangle, rounded corners=3pt, draw=primaryblue!70, fill=primaryblue!8,
             line width=0.8pt, minimum width=26mm, minimum height=8mm,
             font=\scriptsize\bfseries, inner sep=2pt, align=center},
   a/.style={rectangle, rounded corners=2pt, draw=successgreen!70, fill=successgreen!10,
             line width=0.8pt, minimum width=32mm, minimum height=8mm,
             font=\scriptsize, inner sep=2pt, align=center},
   every edge/.style={-Latex, line width=0.8pt, draw=darkgray}]
  \node[q] (root) at (0,0) {Read statement};
  \node[q] (lab)  at (-4,-1.6) {labelled?};
  \node[q] (un)   at (4,-1.6) {given graph\\ \(G\)?};
  \node[a] (pf)  at (-6,-3.4) {Pr\"ufer\\ \(\secref{c32-prufer}\)};
  \node[a] (cy)  at (-2.6,-3.4) {Cayley /\\ completion};
  \node[a] (mt)  at (1.6,-3.4) {Matrix-\\ Tree};
  \node[a] (ex)  at (5.6,-3.4) {Linearity\\ of \(\mathbb{E}\)};
  \node[a] (bo)  at (8,-3.4) {Hand off\\ Ch.28/26/31};
  \draw (root)--(lab)--(pf);
  \draw (lab)--(cy);
  \draw (root)--(un)--(mt);
  \draw (un)--(ex);
  \draw (un.east) to[out=0,in=-90] (bo.south);
\end{tikzpicture}
\end{center}

Every Ch.~30 problem walks one of the four leaf boxes.  The "Hand off"
exit is \(\secref{c32-boundary}\)'s pick list: an unlabelled question marks
a Ch.~28 or Ch.~26 problem dressed as a Ch.~30 one.

\chaptersection{Common Mistakes Checklist}

\begin{enumerate}[nosep]
  \item \textbf{Using Cayley for a \emph{given} graph.}  Cayley \(n^{n-2}\)
        applies only when the graph is \(K_n\); for an arbitrary \(G\) you
        must use \(\thmref{c32-matrix}\) or lose correctness on non-complete
        instances.
  \item \textbf{Dividing a labelled count by \(n!\) to get the unlabelled
        one.}  The orbit sizes are \(n!/|\mathrm{aut}(S)|\), not uniform;
        \(\thmref{c32-orbit}\) has to be summed shape-by-shape.
  \item \textbf{Pr\"ufer code length off by one.}  \(n^{n-2}\) is derived
        from a code of length \(\boldsymbol{n-2}\) over \([n]\); a length
        \(n-1\) code enumerates \emph{rooted} trees.
  \item \textbf{Reading \(P(T)=c_{n-2}\).}  The last-code-entry rule is
        wrong whenever \(c_{n-2}=n\) (see the counterexample in
        \(\secref{c32-cayley}\)).
  \item \textbf{Modular determinant without a prime modulus.}  Gaussian
        elimination requires an invertible pivot; use a prime (\(10^9+7\),
        \(998244353\)) or switch to Bareiss' exact-division variant.
  \item \textbf{Linearity with independence.}  \(\lemref{c32-indexp}\)
        does not assume independence, so a proof that "justifies"
        independence is at best wasted work and at worst hides a bug in
        the decomposition.
  \item \textbf{Depth \(0\) vs.\ depth \(1\) for the root.}  CF 280C's
        expectation is \(\sum_v 1/\mathrm{depth}(v)\) with
        \(\mathrm{depth}(1)=1\); a \(0\)-based DFS depth diverges on the
        root term.
  \item \textbf{Handling \(k=1\) in \(n^{k-2}\).}  The exponent
        becomes \(-1\) at \(k=1\); special-case the "already a tree"
        branch (total \(=1\)).
  \item \textbf{Forgetting the \(s_i\) product in forest-completion.}
        \(n^{k-2}\) alone is the count when each component is a single
        vertex; the product \(\prod s_i\) is the local graft choice.
  \item \textbf{Skipping the binomial inversion in the Summit.}  The DP
        returns \(f(k)\) -- the count of trees \emph{containing} a fixed
        \(k\)-set -- not \(g(k)\), the exact-\(k\) count.  Forgetting the
        inversion produces answers that are systematically too big.
\end{enumerate}

\chaptersection{Wrong-Answer Clinic}

\begin{itemize}[nosep]
  \item \textbf{WA on CF 2190D sample case 1 (n=3, m=0).}  A common cause:
        counting all three labelled trees under the assumption
        \(P(T)=c_{n-2}\).  Codes \([1],[2],[3]\) give \(P(T)=1,2,2\), and
        the "last entry" heuristic assigns \(1,2,3\) -- wrong on the
        third.  The histogram is \((1,2,0)\) instead of the correct
        \((1,2)\) once you exclude \(v=n\).
  \item \textbf{WA on CF 280C sample \(n=2\).}  Emitting \(0.5\) instead
        of \(1.5\): the root \(1\) is at depth \(0\) in a BFS-level view,
        so \(1/\mathrm{depth}\) is either undefined or you have to skip
        the root.  Both are wrong -- use vertex-count depth starting at
        \(1\).
  \item \textbf{WA on the Summit for \(k=0\).}  Reporting \(f(0)=n^{n-2}\)
        as the answer; but \(g(0)\) is not \(f(0)\) -- it is
        \(\sum_j(-1)^j\binom{j}{0}f(j)=\sum_j(-1)^j f(j)\), the
        inclusion-exclusion over all tree edges.  On the star \(n=3\):
        \(f(0)-f(1)+f(2)=3-4+1=0\), matching sample.
  \item \textbf{WA on Matrix-Tree for \(G\) disconnected.}  Your code
        prints the determinant of \(L^{(i)}\) as if it were nonzero; but
        the kernel of a disconnected \(G\)'s Laplacian has dimension \(>1\)
        so every cofactor is zero.  The right answer is \(0\), and any
        nonzero output means you built \(L\) wrong (typically a missing
        degree on the diagonal).
\end{itemize}

\chaptersection{Quick Reference Card}

\begin{quickref}
\begin{itemize}[nosep]
  \item \textbf{Pr\"ufer code.}  Length \(n-2\), labels in \([n]\),
        \(m_v(T)=\deg_T(v)-1\) copies of \(v\).
  \item \textbf{Cayley.}  \(|\mathcal{T}_n|=n^{n-2}\).
  \item \textbf{Rooted.}  \(n^{n-1}\).
  \item \textbf{Degree prescribed.}  \((n-2)!/\prod_i(d_i-1)!\).
  \item \textbf{Forest completion.}  \(n^{k-2}\prod_i s_i\).
  \item \textbf{Complete bipartite \(K_{a,b}\).}  \(a^{b-1}b^{a-1}\).
  \item \textbf{Cycle \(C_n\).}  \(n\) spanning trees.
  \item \textbf{Matrix-Tree.}  \(\tau(G)=\det L^{(i)}(G)\), any \(i\).
  \item \textbf{Linearity.}  \(\mathbb{E}[\sum X_i]=\sum \mathbb{E}[X_i]\),
        no independence needed.
  \item \textbf{CF 280C.}  \(\mathbb{E}=\sum_v 1/\mathrm{depth}(v)\),
        depth in vertices from root.
  \item \textbf{Summit (CF 917D).}  \(f(k)=n^{n-k-2}\sum_S\prod_i s_i\),
        then \(g(k)=\sum_{j\ge k}(-1)^{j-k}\binom{j}{k}f(j)\).
\end{itemize}
\end{quickref}

\chaptersection{Write-From-Memory Drills}

Each drill is a small program you should be able to reproduce on a blank
editor in ten minutes; the labels point back to the code the drill rebuilds.

\begin{enumerate}[nosep]
  \item Encode \([5,5,3,4]\)-style Pr\"ufer code from an edge list --
        rebuild \(\coderef{c32-encode}\) in \(\bigO{n\log n}\) with a
        min-heap of leaves and a degree array.
  \item Decode a Pr\"ufer code back to an edge list -- rebuild
        \(\coderef{c32-decode}\) with a count array and a min-heap of
        zero-count labels; the final two heap entries are \(n\) and \(P(T)\).
  \item Given a Laplacian as a \(2\)D vector, compute the modular
        determinant via Gaussian elimination -- rebuild
        \(\coderef{c32-moddet}\); handle row swaps and pivots explicitly.
  \item Compute \(\sum_v 1/\mathrm{depth}(v)\) on a rooted tree given
        as an edge list -- rebuild \(\coderef{c32-game}\); iterative DFS,
        double accumulator, print with \(20\) digits.
  \item Compute \(n^{k-2}\prod_i s_i\) modulo a prime for a forest --
        rebuild \(\coderef{c32-modpow}\) and its DSU driver from
        \(\coderef{c32-prufvtx}\); special-case \(k=1\).
\end{enumerate}

\chaptersection{Complete Problem Checklist}

The chapter's roster and the exact tool each problem needs; use this as a
one-row-per-problem review before the mixed set.

\begin{center}
\renewcommand{\arraystretch}{1.2}\setlength{\tabcolsep}{3pt}
\begin{tabular}{@{}p{0.16\linewidth} >{\raggedright\arraybackslash}p{0.24\linewidth} >{\raggedright\arraybackslash}p{0.20\linewidth} >{\raggedright\arraybackslash}p{0.10\linewidth} >{\raggedright\arraybackslash}p{0.18\linewidth}@{}}
\toprule
\textbf{ID} & \textbf{Title} & \textbf{Tool} & \textbf{Rating} & \textbf{Section}\\
\midrule
CF 2190D & Prufer Vertex & forest-completion + DP fold & 2500 & 30.2\\
CF 280C  & Game on Tree  & linearity + depth reciprocals & 2200 & 30.4\\
CF 917D  & Stranger Trees & completion + binomial inversion & 2500 & Summit\\
\bottomrule
\end{tabular}
\end{center}

\chaptersection{Recommended Study Order}

Read \(\secref{c32-prufer}\) first and hand-encode the running tree; the
bijection is the one idea the whole chapter rests on.  Then \(\secref
{c32-cayley}\), whose Cayley corollary should feel immediate.  Work the
CF 2190D sample cases by hand before reading the solution -- the
smallest-leaf peel on \(n=3\) and \(n=5\) is quick and pins the notion of
\(P(T)\) far better than a paragraph of prose.

\(\secref{c32-matrix}\) is best approached once \(\secref{c32-cayley}\) has
been digested -- the \(K_n\) sanity check on Cayley is the payoff.  Then
\(\secref{c32-expect}\), which is a short section and a fun change of
register: the same input object (a tree) but the answer is a real number.
Finish with \(\secref{c32-boundary}\), then the Summit.

\chaptersection{Difficulty Progression}

\begin{itemize}[nosep]
  \item \lv{3} Encode/decode the running tree \(T\); hand-compute \(n^{n-2}\)
        for \(n=1..5\) and verify Cayley on the shapes you can draw.
  \item \lv{4} Solve CF 280C in one DFS once you can articulate why
        \(\Pr[A_v]=1/\mathrm{depth}(v)\).  Read Matrix-Tree on \(C_5\),
        verify \(\tau(C_5)=5\) both combinatorially and by hand.
  \item \lv{5} Reconstruct the CF 2190D samples by hand and code the
        forest-completion formula.  Read the Summit's DP structure, and
        verify \(g(0)=0\) on the star \(n=3\) using the inversion.
  \item \lv{6} Solve the Summit CF 917D on your own; hand-code the
        \((k,s)\) DP up to \(n\le 20\) and check \(f(k)\) matches the
        binomial-inversion output.
\end{itemize}

\chaptersection{Reduction Gallery}

\begin{itemize}[nosep]
  \item \emph{Count spanning trees of \(K_n\)} \(\to\) Cayley \(n^{n-2}\)
        \(\to\) Pr\"ufer: sequences of length \(n-2\) over \([n]\).
  \item \emph{Count completions of a forest} \(\to\) contract each
        component to a super-vertex, then Pr\"ufer on the super-tree with
        a local \(s_i\)-way graft per super-vertex.
  \item \emph{Count labelled trees sharing exactly \(k\) edges with a
        given tree \(T\)} \(\to\) inclusion-exclusion / binomial inversion
        on top of "count containing a fixed subset of \(k\) edges"
        \(\to\) forest-completion per subset.
  \item \emph{Expected steps to empty a rooted tree by random subtree
        deletion} \(\to\) sum of indicators \(\mathbf{1}_{A_v}\)
        \(\to\) uniform-permutation symmetry on the root-\(v\) path
        \(\to\) \(1/\mathrm{depth}(v)\) per vertex.
  \item \emph{Count spanning trees of a specific graph \(G\)} \(\to\)
        Laplacian \(L(G)\) \(\to\) delete row/col \(i\) \(\to\)
        modular Gaussian elimination \(\to\) determinant.
\end{itemize}

\chaptersection{Mixed Problem Set}

\begin{mixedset}
\practice{1605D}{Treelabeling}{CF 1605D}{https://codeforces.com/problemset/problem/1605/D}{CF 2100}{greedy, tree DP, label assignment}{Assign letters to tree vertices to minimise the sum along paths; the trick is not Prüfer but a careful greedy on centroid-like structure. Read as a warm-up on tree-DP counting before the harder items.}
\practice{1311E}{Construct the Binary Tree}{CF 1311E}{https://codeforces.com/problemset/problem/1311/E}{CF 2200}{construction, tree DP, level assignment}{Construct a labelled binary tree from a level array; the count is a multinomial over level sizes once you fix the parent-choice per node -- same shape of argument as Prüfer-degree constraints.}
\practice{1540B}{Tree Array}{CF 1540B}{https://codeforces.com/problemset/problem/1540/B}{CF 2300}{tree DP, segment tree, combinatorics}{A tree-DP fold with a segment tree on the Euler tour; cross-cite Ch.19 for the monoid.}
\practice{1929F}{Sasha and the Wedding BST}{CF 1929F}{https://codeforces.com/problemset/problem/1929/F}{CF 2300}{BST count, Catalan-like, tree enumeration}{Count BST labelings on a fixed structure; the multinomial and Catalan-like recurrence are cousins of Cayley's argument.}
\practice{653E}{Bear and Forgotten Tree 2}{CF 653E}{https://codeforces.com/problemset/problem/653/E}{CF 2400}{spanning trees, Matrix-Tree, mod arithmetic}{Direct application of \(\secref{c32-matrix}\): count spanning trees of a given graph modulo a prime via Laplacian cofactor.}
\practice{1237E}{Balanced Binary Search Trees}{CF 1237E}{https://codeforces.com/problemset/problem/1237/E}{CF 2400}{tree counting, DP on sizes}{Count balanced BSTs; recursive size-DP mirrors Cayley's derivation but with a height-balanced filter.}
\practice{1120D}{Power Tree}{CF 1120D}{https://codeforces.com/problemset/problem/1120/D}{CF 2500}{Prüfer-style, MST-style greedy, tree reconstruction}{Given a distance matrix, reconstruct the tree with minimum weight; use Prüfer-like degree reasoning plus a Kruskal-style filter. Cite Ch.31.}
\practice{804D}{Expected diameter of a tree}{CF 804D}{https://codeforces.com/problemset/problem/804/D}{CF 2500}{linearity of expectation, edge indicators}{Split diameter into edge indicators; the technique is exactly \(\secref{c32-expect}\) but with edge indicators on path-uses.}
\practice{1762E}{Tree Sum}{CF 1762E}{https://codeforces.com/problemset/problem/1762/E}{CF 2600}{subtree reweight DP, tree DP}{Reweight subtree contributions and recompute the maximum path sum; a fold-heavy Ch.13 style problem with an enumerative flavour.}
\practice{1495D}{BFS Trees}{CF 1495D}{https://codeforces.com/problemset/problem/1495/D}{CF 2600}{BFS tree counting, combinatorics}{Count BFS-tree realizations of a given graph level structure; the per-level product mirrors the forest-completion \(\prod s_i\) factor.}
\end{mixedset}

\chaptersection{The Summit -- CF 917D Stranger Trees}

\begin{notebox}[Why this is the chapter top rope]
CF 917D asks a counting question whose surface is unlabelled ("share
exactly \(k\) edges with a given tree") but whose answer space is
\emph{labelled trees on \([n]\)}.  The tool is \(\thmref{c32-forest}\) (the
forest-completion formula from \(\secref{c32-cayley}\)) applied per subset
of the given tree's edges, combined with a \((k,s)\)\emph{-shaped} tree DP
that folds the product \(\prod_i s_i\) across subsets, plus a binomial
inversion to move from "at least" to "exactly".  Three tools from this
chapter (the bijection's completion corollary, a tree DP, and modular
arithmetic) and one from Ch.~13 (the bottom-up fold \(\chref{c10}\)) make
the full pipeline.
\end{notebox}

\solvedproblem{c32-summit}{Stranger Trees}{CF 917D}{https://codeforces.com/problemset/problem/917/D}{CF 2500 L5}{Pruefer bijection, forest-completion count \(\thmref{c32-forest}\), tree DP over subsets, binomial inversion}

\begin{problemstatement}[CF 917D, verified against the live CF problem page]
A labelled tree \(T\) is given on \(n\) vertices with \(1\le n\le 100\); the
next \(n-1\) lines describe its edges as pairs \((v,u)\) with \(1\le v,u\le
n\), \(v\ne u\).  For each \(k\in[0,n-1]\), count the number of
\emph{labelled} trees \(T'\in\mathcal{T}_n\) that share \emph{exactly} \(k\)
edges with \(T\) -- two labelled trees are considered different iff their
edge sets differ.  Print the \(n\) answers modulo \(10^9+7\).  Time limit
1 second, memory 256 MB.  Samples: (i) star \(1\!-\!2,\, 1\!-\!3\) on
\(n=3\) prints \(0\;2\;1\); (ii) path \(1\!-\!2,\, 2\!-\!3,\, 3\!-\!4\) on
\(n=4\) prints \(1\;7\;7\;1\); (iii) star \(1\!-\!2,\, 1\!-\!3,\, 1\!-\!4\)
on \(n=4\) prints \(0\;9\;6\;1\).
\end{problemstatement}

\begin{figure}[H]
\centering
\begin{tikzpicture}[diagram, scale=0.9, transform shape,
   tnode/.style={circle, draw=primaryblue!70, fill=primaryblue!8, line width=0.7pt,
                 minimum size=6.5mm, font=\scriptsize\bfseries, inner sep=1pt}]
  \node[tnode] (c) at (0,0) {1};
  \node[tnode] (l1) at (-1.4,1.2) {2};
  \node[tnode] (l2) at (1.4,1.2) {3};
  \draw[primaryblue!70, line width=1pt] (c)--(l1) (c)--(l2);
  \node[annot, font=\scriptsize] at (0,-1.2) {star on 3: answers 0 2 1};
  \node[tnode] (p1) at (4,1.2) {1};
  \node[tnode] (p2) at (5,0.4) {2};
  \node[tnode] (p3) at (6,-0.4) {3};
  \node[tnode] (p4) at (7,-1.2) {4};
  \draw[primaryblue!70, line width=1pt] (p1)--(p2)--(p3)--(p4);
  \node[annot, font=\scriptsize] at (5.5,-2.1) {path on 4: answers 1 7 7 1};
\end{tikzpicture}
\caption{Two of the three official samples; the third is the four-leaf star,
whose histogram is \(0\;9\;6\;1\).  The Summit recovers all three via the Solution step's pipeline.}
\label{fig:c32-summit}
\end{figure}

\begin{attemptgate}[Attempt it yourself first -- \lv{5}, give yourself 60 minutes]
Before coding: what quantity of the forest \((V,S)\) for a subset \(S\subseteq
E(T)\) controls the number of completions that contain \(S\)?  And how do
you convert "count completions containing \(S\)" into "count completions
sharing \emph{exactly} \(k\) edges with \(T\)"?
\peek{\(\thmref{c32-forest}\) gives the count of trees containing a fixed
forest \((V,S)\) directly -- the answer depends on the sizes of \((V,S)\)'s
components, not their shape.  For the second question: binomial inversion
between "at least" and "exactly" counts.}
\nudge{Define \(f(k)=\sum_{S\subseteq E(T),|S|=k}\) (trees containing \(S\)).
Define \(g(j)=\) the number of
trees sharing exactly \(j\) edges with \(T\).  Then
\(f(k)=\sum_{j\ge k}\binom{j}{k}g(j)\), because each tree sharing \(j\) edges
contributes \(\binom{j}{k}\) subsets of size \(k\).  Invert.
}
\fullhint{For \(f(k)\): iterate over \(S\subseteq E(T)\) of size \(k\),
which is a forest with \(n-k\) components of sizes \(s_1,\dots,s_{n-k}\);
\(\thmref{c32-forest}\) yields \(n^{n-k-2}\prod_i s_i\) completions per
subset.  Aggregate over all \(S\) of size \(k\) with a tree DP: root \(T\)
at \(1\), and let \(dp[u][j][t]\) be the sum over size-\(j\) subsets of
edges inside \(u\)'s subtree of the product of closed component sizes,
with the component containing \(u\) currently having size \(t\) (still
open).  Merge child subtrees by either dropping the connecting edge (the
child's open becomes closed, multiply by its open size) or keeping it
(the opens merge, \(t\) adds).  Final answer: \(f(k)=n^{n-k-2}\sum_t t\cdot
dp[1][k][t]\).  Then invert binomially.}
\end{attemptgate}

\solstep{Understanding}
The Summit is a binomial inversion wearing a Pr"ufer hat.  The natural
quantity to enumerate is \(f(k)\): the total over all \(k\)-subsets \(S\subseteq
E(T)\) of the number of labelled trees that \emph{contain} \(S\).  Each
such tree \(T'\) sharing exactly \(j\) edges with \(T\) contributes to
\(f(k)\) once per \(k\)-subset of its shared \(j\)-edge intersection, i.e.\
\(\binom{j}{k}\) times.  Summing gives \(f(k)=\sum_{j\ge k}\binom{j}{k}
g(j)\), and the standard binomial inversion recovers \(g(k)=\sum_{j\ge k
}(-1)^{j-k}\binom{j}{k}f(j)\).  The problem reduces to computing \(f(k)\)
for \(k=0,\dots,n-1\), which \(\thmref{c32-forest}\) makes a subtree-DP
exercise on \(T\) itself.

\solstep{Key Observations}
\begin{observation}
\(f(k)\) depends on \(S\) only through the \emph{component sizes} of the
forest \((V,S)\): the completion count \(n^{n-k-2}\prod_i s_i\) is
invariant under tree-shape changes that preserve the size multiset.
\end{observation}
\begin{observation}
The DP \(dp[u][j][t]\) has \(O(n^3)\) states per node (bounded by
subtree size), and the merge over a child's DP is \(O(n^2)\), so total
time is \(O(n^3)\), which fits \(n\le 100\) with room.
\end{observation}
\begin{observation}
Binomial inversion uses only \(\binom{n}{k}\bmod p\) with \(p=10^9+7\); precompute
factorials and inverse factorials in \(O(n)\) and query in \(O(1)\).
\end{observation}

\solstep{Think It With Me}
\begin{thinkbox}
\deadend{Enumerate all \(S\subseteq E(T)\)}{\(2^{n-1}\) subsets is \(2^{99}\)
for \(n=100\), hopeless.}
\deadend{Compute \(g(k)\) directly via a filter DP}{The "exactly \(k\)\
" filter requires tracking which tree edges are shared with \(T'\), but
\(T'\) is not explicit -- only its overlap with \(T\) matters, and there's
no way to enumerate \(T'\)'s edges without leaving labelled-tree space.}
\deadend{Use Matrix-Tree on a modified graph}{Matrix-Tree counts spanning
trees of a fixed \(G\), which does not match the question "count \(T'\)\
sharing exactly \(k\) edges with \(T\)" -- you would need to iterate over
exponentially many \(G\)'s.}
\aha{\(f(k)\) is a tree DP on \(T\) itself, with state \((j,t)\) per node
and a merge that either closes the child's open component or fuses it into
the parent's open.  Then binomial inversion converts \(f\) to \(g\).}
\end{thinkbox}

\solstep{The Approach Ladder}
\begin{center}\renewcommand{\arraystretch}{1.25}\setlength{\tabcolsep}{3pt}
\begin{tabular}{@{}c p{0.26\linewidth} p{0.30\linewidth} l c@{}}
\toprule
\textbf{\#} & \textbf{Approach} & \textbf{Idea} & \textbf{Time} & \textbf{Verdict}\\
\midrule
1 & \lv{2}\ enumerate all \(T'\) & \(n^{n-2}\) candidates & \(\bigO{n^{n}}\) & \tle\\
2 & \lv{3}\ enumerate all subsets \(S\subseteq E(T)\) & \(2^{n-1}\) subsets & \(\bigO{2^n}\) & \tle\\
3 & \lv{4}\ Matrix-Tree on every subgraph & exponential \(G\)'s & \(\bigO{2^n n^3}\) & \tle\\
4 & \lv{5}\ \textbf{tree DP + binomial inversion} & \((j,t)\) state per node & \(\bigO{n^3}\) & \ac\\
\bottomrule
\end{tabular}\end{center}
Rungs 1--3 all blow up on the exponential state space.  Rung 4 is a
tree-DP fold and an \(O(n^2)\) inversion, both \(\bigO{n^3}\) end to end.

\solstep{Solution}
Root \(T\) at vertex \(1\).  Compute the DP \(dp[u][j][t]\) as described.
\(f(k)=n^{n-k-2}\sum_t t\cdot dp[1][k][t]\bmod p\).  Precompute \(\binom{n}{k}\)
and \((-1)^k\) tables.  Invert binomially: \(g(k)=\sum_{j\ge k}(-1)^{j-k
}\binom{j}{k}f(j)\).  Print \(g(0),\dots,g(n-1)\).

\solstep{Implementation}
\begin{lstlisting}[caption={CF 917D. Tree DP with state (edges kept, open size) then binomial inversion. Handles the k=n-1 case by special-casing the negative exponent via modPow with MOD minus 2. ASCII only, identifiers without underscores}, label={lst:c32-summit}]
#include <bits/stdc++.h>
using namespace std;
typedef long long ll;
const ll MOD=1000000007;
ll modPow(ll b, ll e){ ll r=1; b%=MOD; while(e>0){ if(e&1) r=r*b%MOD; b=b*b%MOD; e>>=1;} return r; }
int n;
vector<vector<int> > adj;
vector<vector<vector<ll> > > dp; // dp[u][j][t]
vector<int> subSz;
void dfs(int u, int p){
  subSz[u]=1;
  dp[u][0][1]=1;
  int used=0;
  for(int v : adj[u]) if(v!=p){
    dfs(v,u);
    vector<vector<vector<ll> > > nDp(n+1, vector<vector<ll> >(n+1, vector<ll>(n+1,0)));
    for(int j1=0; j1<=used; j1++)
      for(int t1=1; t1<=subSz[u]; t1++) if(dp[u][j1][t1])
        for(int j2=0; j2<=subSz[v]-1; j2++)
          for(int t2=1; t2<=subSz[v]; t2++) if(dp[v][j2][t2]){
            ll ways=dp[u][j1][t1]*dp[v][j2][t2]%MOD;
            // drop edge (u,v): v's open component becomes closed -> * t2
            nDp[u][j1+j2][t1] = (nDp[u][j1+j2][t1] + ways*t2)%MOD;
            // keep edge (u,v): opens merge
            nDp[u][j1+j2+1][t1+t2] = (nDp[u][j1+j2+1][t1+t2] + ways)%MOD;
          }
    for(int j=0;j<=n;j++) for(int t=0;t<=n;t++) dp[u][j][t]=nDp[u][j][t];
    used += subSz[v];
    subSz[u]+=subSz[v];
  }
}
int main(){
  if(scanf("%d",&n)!=1) return 0;
  adj.assign(n+1, {});
  for(int i=0;i<n-1;i++){ int u,v; scanf("%d %d",&u,&v); adj[u].push_back(v); adj[v].push_back(u); }
  dp.assign(n+1, vector<vector<ll> >(n+1, vector<ll>(n+1,0)));
  subSz.assign(n+1,0);
  dfs(1,0);
  vector<ll> f(n+1,0);
  for(int k=0;k<=n-1;k++){
    ll sum=0;
    for(int t=1;t<=n;t++) sum=(sum + (ll)t*dp[1][k][t])%MOD;
    ll c=(k==n-1)?1:modPow(n, n-k-2);
    f[k]=sum*c%MOD;
  }
  // binomial inversion g(k) = sum_{j>=k} (-1)^{j-k} C(j,k) f(j)
  vector<ll> fact(n+1,1), invFact(n+1,1);
  for(int i=1;i<=n;i++) fact[i]=fact[i-1]*i%MOD;
  invFact[n]=modPow(fact[n], MOD-2);
  for(int i=n;i>=1;i--) invFact[i-1]=invFact[i]*i%MOD;
  auto nCk=[&](int a,int b)->ll{ if(b<0||b>a) return 0; return fact[a]*invFact[b]%MOD*invFact[a-b]%MOD; };
  for(int k=0;k<=n-1;k++){
    ll gk=0;
    for(int j=k;j<=n-1;j++){
      ll term=nCk(j,k)*f[j]%MOD;
      if((j-k)&1) gk=(gk-term+MOD)%MOD; else gk=(gk+term)%MOD;
    }
    printf("%lld ", gk);
  }
  printf("\n");
  return 0;
}
\end{lstlisting}

\solstep{Trace on the Official Samples}
\begin{tracebox}[Sample 1: \(n=3\) star \(1\!-\!2,\,1\!-\!3\)]
DP gives \(f(0)=3^{1}\cdot 1\cdot 1\cdot 1=3\) (empty subset; three isolates;
\(\prod s_i=1\)).  \(f(1)\) sums over two 1-edge subsets, each giving
components \((2,1)\) with \(\prod s_i=2\), times \(3^{0}=1\): \(f(1)=2+2=4\).
\(f(2)=1\) (the whole star; \(c=1\); the exponent \(-1\) is special-cased
with \(\prod s_i=3\), giving \(3\cdot 3^{-1}=1\)).  Inversion: \(g(0)=f(0)-f(1
)+f(2)=3-4+1=0\); \(g(1)=f(1)-2f(2)=4-2=2\); \(g(2)=f(2)=1\).  Output
\(0\;2\;1\), matching the sample.
\end{tracebox}
\begin{tracebox}[Sample 2: \(n=4\) path \(1\!-\!2,\,2\!-\!3,\,3\!-\!4\)]
\(f(0)=4^{2}=16\); \(f(1)=3\cdot 8=24\) (three 1-edge subsets, each \(c=3\), \(\prod s_i=2\), times \(4^{1}=4\), giving \(8\) per subset); \(f(2)=
3+3+4=10\) (two adjacent-edge subsets give \(c=2,\prod s_i=3\) each; one
disjoint subset gives \(\prod s_i=4\)); \(f(3)=1\).  Inversion: \(g(0
)=16-24+10-1=1\); \(g(1)=24-20+3=7\); \(g(2)=10-3=7\); \(g(3)=1\).  Output \(1\;7\;7\;1\),
matching the sample.
\end{tracebox}
\begin{tracebox}[Sample 3: \(n=4\) star \(1\!-\!2,\,1\!-\!3,\,1\!-\!4\)]
\(f(0)=16\); \(f(1)=24\) (three 1-edge subsets, same as Sample 2); \(f(2
)=9\) (all three 2-edge subsets are the two-leaf star, \(c=2\) and \(\prod s_i=3\) each); \(f(3)=1\).  Inversion: \(g(0)=16-24+9-1=0\); \(g(1)=24-18+3=9\); \(g(2)=9-3=6\); \(g(3)=1\).  Output
\(0\;9\;6\;1\), matching the sample.
\end{tracebox}

\solstep{Complexity}
\begin{complexitybox}
The DFS visits \(O(n)\) nodes; the child merge at node \(u\) is \(O((s_u)\cdot s_v)\) states on each side, so the total across the tree is
\(O(n^3)\).  Binomial inversion is \(O(n^2)\).  Space \(O(n^3)\) for the
DP array, \(O(n)\) for the precomputed factorials.  For \(n=100\) the
DP is \(\sim 10^6\) states, comfortably inside 1 second.  All arithmetic
mod \(10^9+7\); use 64-bit throughout.
\end{complexitybox}

\begin{hardeningbox}
\begin{itemize}[nosep,leftmargin=1.4em]
\item \textbf{The \(k=n-1\) exponent.}  \(n^{n-k-2}\) reads as \(n^{-1}\)
when \(k=n-1\); special-case to \(1\) (already a tree, one completion).
\item \textbf{Reset the DP.}  \(dp[u][j][t]\) is per-node; allocate
\((n+1)^3\) once and initialise to zero, or use a per-node vector.
\item \textbf{Child loop bound.}  \(j_2\) runs over \([0, s_v-1]\), not
\([0,n]\); tightening this bound is the difference between \(\bigO{n^4}\)
and \(\bigO{n^3}\).
\item \textbf{Sign on the inversion.}  \((-1)^{j-k}\) is a mod-subtraction;
always write \((x-y+\text{MOD})\%\text{MOD}\) to keep the value nonnegative.
\item \textbf{Rooted at 1.}  The DP assumes a root; picking a different
root permutes which open size at \(u\) corresponds to which component,
but the final sum over \(t\cdot dp[1][k][t]\) is invariant -- still,
fix a root for the code.
\end{itemize}
\end{hardeningbox}

\begin{recallbox}[Chapter 30 in ten seconds]
\emph{Pr"ufer is the bijection; Cayley is the corollary; the forest-
completion formula is the workhorse; Matrix-Tree hands you a determinant
when the graph is fixed; linearity of expectation turns an exponential
state process into a sum of reciprocals of depths.  CF 2190D reads
Pr"ufer-style per vertex, CF 280C is one DFS on depth reciprocals, and
the Summit CF 917D is a tree DP with a binomial inversion.}  The
chapter ends with the same object it started with -- a labelled tree on
\([n]\) -- but you can now \emph{count} them under any reasonable filter
and \emph{expect} things about them without writing a recursion.
\selftest{
  \item State the completion formula \(\thmref{c32-forest}\) and its
        \(k=n\) and \(k=1\) special cases.
  \item Give the binomial-inversion identity relating \(f(k)\) and \(g(k)\).
  \item Explain why a naive \(P(T)=c_{n-2}\) filter fails on CF 2190D.
  \item Reconstruct the running example \(n=3\) star Summit pipeline
        from memory.
}
\end{recallbox}

\chaptersection{Tool Kit Inventory}

\begin{center}
\renewcommand{\arraystretch}{1.2}\setlength{\tabcolsep}{3pt}
\begin{tabular}{@{}p{0.28\linewidth} >{\raggedright\arraybackslash}p{0.22\linewidth} >{\raggedright\arraybackslash}p{0.36\linewidth}@{}}
\toprule
\textbf{Tool} & \textbf{Owned?} & \textbf{Use in this chapter}\\
\midrule
Pr"ufer encode/decode & Yes, \(\secref{c32-prufer}\) & the whole chapter's frame \\
Cayley \(n^{n-2}\) & Yes, \(\thmref{c32-cayley}\) & base count \\
Rooted \(n^{n-1}\) & Yes, \(\thmref{c32-rooted}\) & normaliser \\
Degree multinomial & Yes, \(\corref{c32-mult}\) & CF 2190D filter \\
Forest completion & Yes, \(\thmref{c32-forest}\) & CF 2190D, Summit \\
Laplacian + Matrix-Tree & Yes, \(\secref{c32-matrix}\) & given-graph counts \\
Modular determinant & Yes, \(\coderef{c32-moddet}\) & Matrix-Tree engine \\
Linearity of expectation & Yes, \(\lemref{c32-indexp}\) & CF 280C \\
Binomial inversion & Yes, Summit & \(f\to g\) filter \\
Depth from root & no, Ch.12 (\chref{c09}) & \(\secref{c32-expect}\) reads it \\
Tree DP fold & no, Ch.13 (\chref{c10}) & Summit merges via it \\
DSU component census & no, Ch.17 (\chref{c28}) & CF 2190D forest read \\
Segment tree monoid & no, Ch.19 (\chref{c13}) & cited once, not required here \\
AHU tree isomorphism & no, Ch.28 (\chref{c22}) & \(\secref{c32-boundary}\) hand-off \\
\bottomrule
\end{tabular}
\end{center}

\chaptersection{Related Topics}

\begin{itemize}[nosep]
  \item Ch.~13 (\(\chref{c10}\)) tree DP -- the bottom-up fold; the Summit's
        \((j,t)\) DP is that fold with an "open component size" state
        layered on top.
  \item Ch.~17 (\(\chref{c28}\)) DSU and small-to-large -- CF 2190D reads
        the forest's components via DSU; the fold's merge can also be
        sped up via small-to-large.
  \item Ch.~26 (\(\chref{c30}\)) virtual trees -- the handoff in
        \(\secref{c32-boundary}\) when the query is over a marked vertex
        subset rather than the whole \([n]\).
  \item Ch.~28 (\(\chref{c22}\)) tree isomorphism and AHU -- the unlabelled
        side of \(\thmref{c32-orbit}\); needed to compute \(|\mathrm{aut}(S)|\) when reconciling labelled and unlabelled counts.
  \item Ch.~29 (\(\chref{c31}\)) trees from graphs -- the previous chapter;
        it ends where Ch.~30 starts (a tree is given); Ch.~30 begins at
        the count.
  \item Ch.~31 (\(\chref{c17}\)) MST and variants -- the next chapter;
        a different "tree inside a graph" that counts via a greedy, not
        a determinant.
  \item Ch.~12 (\(\chref{c09}\)) tree metrics -- depth array used verbatim
        by \(\secref{c32-expect}\).
  \item Ch.~19 (\(\chref{c13}\)) segment trees -- cited for the range-max
        monoid used in Chapter 29's Summit; no segment tree is needed
        inside Chapter 30 itself.
\end{itemize}

\begin{lastword}[Last Word --- Chapter 30]
Chapter~30 closes the enumerative loop opened in Chapter~1: the same
labelled tree you first met as an input object is here a \emph{counting
object}, and you own four ways to answer questions about the family.  The
Pr"ufer bijection is the deepest; every other tool in this chapter is
either a corollary of it (Cayley, the completion formula) or a sibling
engine for the case it does not cover (Matrix-Tree for a given graph,
linearity for an expectation).  The Summit CF 917D is a deliberate
capstone: it uses three of the four tools in one pipeline, and the fourth
(matrix-tree) is what you would reach for if the input were not a \emph{
tree} but a graph.  The next chapter, Ch.~31 (\(\chref{c17}\)), turns the
"given graph" branch into a weighted one and hands you Kruskal, Prim, and
the reconstruction tree -- the optimisation twin of the enumeration you
have just finished.
\end{lastword}



\closechapter
